% concatenated from maindegree2.tex
\documentclass[twoside]{amsart}

% ----------------- Pacotes básicos -----------------
\usepackage[utf8]{inputenc}
\usepackage[T1]{fontenc}
\usepackage{t1enc}


% Math
\usepackage{amsmath,amssymb,amsfonts,amsthm,amscd,mathrsfs,mathtools}
\numberwithin{equation}{section}

% Letras caligráficas e afins
\usepackage[mathscr]{eucal}

% Layout e texto
\usepackage{indentfirst}
\usepackage{relsize}
\usepackage[tablename=Table]{caption}
\usepackage[active]{srcltx}

% Gráficos e figuras
\usepackage{graphicx}
\usepackage{graphics}
\usepackage{pict2e}
\usepackage{epic}
\usepackage{epstopdf}
\usepackage{tabularx}
\usepackage{booktabs,makecell,siunitx,eqparbox}

% Cores
\usepackage[x11names,dvipsnames]{xcolor}


% Diagramas
\usepackage[all]{xy}
\usepackage{tikz}
\usepackage{tikz-cd}
\usetikzlibrary{
  matrix,
  calc,
  arrows,
  arrows.meta,
  positioning,
  shapes,
  decorations.pathmorphing
}

% Showkeys
%\usepackage[notref,notcite]{showkeys}

% Ifpdf (sem hiperref aqui ainda)
\usepackage{ifpdf}

%---- listas
\usepackage{enumitem}
\setlist[enumerate]{label=\arabic*)}


% ----------------- Macros simples -----------------
\def\wh{\widehat}
\newcommand{\ba}{\begin{array}}
\newcommand{\ea}{\end{array}}

\newcommand{\tabhead}{\textbf}

% Corrigido: uma barra só em \mathbb
\newcommand{\numset}[1]{\mathbb{#1}}

% ----------------- Seus atalhos geométricos -----------------
\newcommand{\op}[1]{{\mathcal O}_{\mathbb{P}^{#1}}}
\newcommand{\p}[1]{{\mathbb{P}^{#1}}}
\newcommand{\opn}{{\mathcal O}_{\mathbb{P}^{n}}}
\newcommand{\pn}{{\mathbb{P}^{n}}}
\newcommand{\tp}[1]{{\rm T}\mathbb{P}^{#1}}
\newcommand{\tpn}{T{\mathbb{P}^{n}}}
\newcommand{\om}[2]{\Omega_{\mathbb{P}^{#2}}^{#1}}
\newcommand{\omn}[1]{\Omega_{\mathbb{P}^{n}}^{#1}}

\newcommand{\sI}{\mathscr{I}}
\newcommand{\iz}{\mathscr{I}_{Z}}
\newcommand{\ic}{\mathscr{I}_{C}}
\newcommand{\quotf}{{\mathcal Q}{\it uot}}
\newcommand{\dist}{{\mathcal D}{\it ist}}
\newcommand{\sch}{\mathfrak{Sch}_{/\mathbb{C}}}

\DeclareMathOperator{\sing}{Sing}
\DeclareMathOperator{\Supp}{Supp}
\DeclareMathOperator{\im}{im}
\DeclareMathOperator{\rk}{rk}
\DeclareMathOperator{\Coker}{Coker}
\DeclareMathOperator{\Hom}{Hom}
\DeclareMathOperator{\ch}{ch} 
\DeclareMathOperator{\td}{td}
\DeclareMathOperator{\Suporte}{Supp}
\DeclareMathOperator{\codim}{codim}
\DeclareMathOperator{\Ext}{Ext}
\DeclareMathOperator{\Spec}{Spec}
\DeclareMathOperator{\Sym}{Sym}
\DeclareMathOperator{\id}{id}
\DeclareMathOperator{\Res}{Res}
\DeclareMathOperator{\Reg}{Reg}
\DeclareMathOperator{\norm}{norm}
\DeclareMathOperator{\coker}{coker}
\DeclareMathOperator{\rad}{rad}
\DeclareMathOperator{\pic}{Pic}
\DeclareMathOperator{\Div}{Div}
\DeclareMathOperator{\depth}{depth}
\DeclareMathOperator{\Tor}{Tor}
\DeclareMathOperator{\Var}{Var}
\DeclareMathOperator{\Ann}{Ann}
\DeclareMathOperator{\End}{End}

\newcommand{\Hilb}{\operatorname{Hilb}}
\newcommand{\inhom}{{\mathcal H}{\it om}}
\newcommand{\inext}{{\mathcal E}{\it xt}}

\newcommand{\sln}{\mathfrak{sl}}
\newcommand{\mf}[1]{\mathfrak{#1}}

\def\PP{{\mathbb{P}}}
\def\OO{\mathcal{O}}
\def\CC{{\mathbb{C}}}
\def\RR{{\mathbb{R}}}
\def\ZZ{{\mathbb{Z}}}
\newcommand{\QQ}{\mathbb{Q}}

\newcommand{\sD}{\mathscr{D}}
\newcommand{\sF}{\mathscr{F}}
\newcommand{\sG}{\mathscr{G}}
\newcommand{\sH}{\mathscr{H}}
\newcommand{\tD}{{\rm T}\sD}
\newcommand{\tF}{{\rm T}\sF}
\newcommand{\tG}{{\rm T}\sG}
\newcommand{\tH}{{\rm T}\sH}
\newcommand{\nF}{{\rm N}\sF}
\newcommand{\nG}{{\rm N}\sG}
\newcommand{\nH}{{\rm N}\sH}

\newcommand{\Fol}{\mathbb{F}{\rm ol}}
\newcommand{\Log}{\mathcal{L}{\it og}}
\newcommand{\SL}{\mathscr{L}{\rm Lie}}
\newcommand{\AL}{\mathcal{AL}}

\newcommand{\tX}{{\rm T}X}

\setcounter{tocdepth}{1}

\newcommand{\into}{\hookrightarrow}
\newcommand{\onto}{\twoheadrightarrow}
\newcommand*{\ponto}{\makebox[1.5ex]{\textbf{$\cdot$}}}
\newcommand{\diam}{{\rm{diam}}}
\newcommand{\ovl}{\overline}

\newcommand{\ld}{\mathcal L}
\newcommand{\lra}{ \longrightarrow }
\newcommand{\aut}{\mathfrak{aut}}
\newcommand{\fix}{\mathfrak{fix}}
\newcommand{\afix}{\mathfrak{afix}}
\DeclareMathOperator{\Aut}{Aut}
\DeclareMathOperator{\PGL}{PGL}
\newcommand{\lieder}{\mathcal{L}}
\newcommand{\extp}{{\textstyle \bigwedge}} 
\newcommand{\extd}{{\textstyle \bigwedge}^{[2]}} 
\newcommand{\del}[1]{\frac{\partial}{\partial #1}}

\newcommand{\lpb}[3]{{\rm LPB }(#1,#2,#3)}
\newcommand{\rat}[2]{{\rm Rat}(#1;#2)}
\newcommand{\logg}[3]{{\rm Log}(#1;#2,#3)}
\newcommand{\exc}[3]{{\rm ALie}(#1,#2;#3)}

\newcommand{\red}[1]{\textbf{\color{red} #1 }}

\setlength{\textwidth}{15cm}
\setlength{\textheight}{20.5cm}
\setlength{\oddsidemargin}{0.25in}
\setlength{\evensidemargin}{0.25in}

% ----------------- Estilos TikZ globais -----------------
\tikzset{
  >=latex,
  vertex/.style={circle,draw,inner sep=2pt},
  active/.style={thick},
  inactive/.style={dashed}
}


% estilos globais para esta figura
\tikzset{
  vertex/.style={circle,draw,inner sep=1pt,fill=white},
  edgeActive/.style={line width=0.9pt,color=RoyalBlue3},
  edgeInactive/.style={line width=0.4pt,color=Gray},
  every node/.style={font=\scriptsize}
}

% ----------------- Teoremas -----------------
\newtheorem{theorem}{Theorem}[section]
\newtheorem{proposition}[theorem]{Proposition}
\newtheorem{corollary}[theorem]{Corollary}
\newtheorem{lemma}[theorem]{Lemma}

\newtheorem{innerthmintro}{Theorem}
\newenvironment{thmintro}[1]
{\renewcommand\theinnerthmintro{#1}\innerthmintro}
{\endinnerthmintro}

\newtheorem{innercorintro}{Corollary}
\newenvironment{corintro}[1]
{\renewcommand\theinnercorintro{#1}\innercorintro}
{\endinnercorintro}

\newtheorem{innerpropintro}{Proposition}
\newenvironment{propintro}[1]
{\renewcommand\theinnerpropintro{#1}\innerpropintro}
{\endinnerpropintro}

\theoremstyle{definition}
\newtheorem{definition}{Definition}[section]
\newtheorem{conjecture}{Conjecture}
\newtheorem{example}{Example}[section]
\newtheorem{remark}{Remark}[section]

% ----------------- Hyperref (apenas uma vez) -----------------
\usepackage{hyperref}
\hypersetup{
  colorlinks=true,
  linkcolor=blue,   % cor dos links internos
  filecolor=magenta,
  urlcolor=cyan
}

\date{}


\begin{document}

\title{ Holomorphic foliations of degree two and arbitrary dimension }


\author[M. Corr\^ea]{ Maur\'icio Corr\^ea }
\address[M. Corr\^ea]{Dipartimento di Matematica, Universit\`a degli Studi di Bari,  Via E. Orabona 4, I-70125, Bari, Italy
}
\email{mauricio.correa.mat@gmail.com, mauricio.barros@uniba.it} 

\author[A. Muniz]{Alan Muniz }
\address[A. Muniz]{Universidade Estadual de Campinas (UNICAMP) \\ Instituto de Matemática, Estatística e Computação Científica (IMECC) \\ Departamento de Matem\'atica \\
Rua S\'ergio Buarque de Holanda, 651\\ 13083-859 Campinas-SP, Brasil}
\address{Departamento de Matem\'atica \\ Centro de Ci\^encias Exatas e da Natureza \\ Universidade Federal de Pernambuco \\ Recife - PE, CEP 50740-560, Brasil }
\email{alannmuniz@gmail.com, alan.nmuniz@ufpe.br }
 




\subjclass[2020]{Primary: 37F75,
 58A17,14D20,14J60. Secondary:53D17, 14F06}


\keywords{ Holomorphic foliations, Poisson structures, Moduli spaces} 


\begin{abstract} 
We prove a complete classification for degree-$2$ foliations on $\mathbb{P}^n$ of any dimension, provided they are not algebraically integrable. If $\sF$ is such a foliation, then either $\sF$ is the linear pull-back of a degree 2  foliation by curves on $\mathbb{P}^{n-k+1}$, or a logarithmic foliation of type $(1^{n-k +1},2)$, or a logarithmic foliation of type $(1^{n-k+3})$,  or a linear pull-back of a degree-2 foliation of dimension 2 on $\mathbb{P}^{n-k+2}$ tangent to an action of the Lie algebra $\mathfrak{aff}(\CC)$. Meanwhile, we prove that any 2-dimensional foliation tangent to a global vector field must satisfy that its tangent sheaf either is not locally free or has a direct summand isomorphic to $\opn(a)$, $a\in\{0,1\}$. As a byproduct of our classification, we describe the geometry of Poisson structures on $\mathbb{P}^{4}$ with generic rank two.
\end{abstract}

\maketitle



\tableofcontents


\section{Introduction}
One of the central problems in the global theory of holomorphic foliations is the description and classification of the irreducible components of the space of foliations on the projective space $\pn$, $n \geq 3$, at least for low degrees. A complete classification is known for degrees $0$ and $1$ in any dimension, and for degree $2$ only in dimension $n-1$. The space $\Fol(d,k,n)$ of degree-$d$ dimension-$k$ foliations on $\pn$ is a projective subvariety of $\PP H^0(\Omega_\pn^p(n-k+d+1))$ whose general element is an integrable $(n-k)$-form without zeros of codimension one. A classical result of Darboux implies that every degree-$0$ foliation is given by the fibers of a linear rational map $\pn \dashrightarrow \p{n-p}$, thus $\Fol(0,k,n)$ is isomorphic to the Grassmannian $G(k,n+1)$ in its Pl\"ucker embedding, see \cite[ Th\'eor\`eme 3.8]{CD2}. In degree $d=1$, Loray, Pereira, and Touzet \cite{LPT2}, generalizing upon Jouanolou \cite{Jou}, proved that $\Fol(1,k,n)$ has two irreducible components: one parameterizing linear pullbacks of degree-one foliations by curves on $\p{n-k+1}$; the other parameterizing rational foliations of type $(1^{n-k-1},2)$, given by the fibers of a rational map $\pn \dashrightarrow \PP(1^{n-k-1},2)$. For $d=2$ and $p=1$, Cerveau and Lins Neto \cite{CLN} proved that $\Fol(2,n-1,n)$ has six irreducible components and described their general element. Since then, several works investigated the classification in dimension $n-1$ and higher degrees \cite{CLNdeg3,  AVdeg4, LPT3}; most remarkably \cite{CLPdeg3} where the authors present an almost complete classification for degree $3$. In a different direction, the present work focuses on degree-$2$ foliations of dimension $k \leq n-2$. The main result of this work is the following.

\begin{thmintro}{A}\label{thm:mainA}
Let $\sF$ be a degree-$2$ foliation on $\pn$ of dimension $k\geq2$. Then one of the following holds:
\begin{enumerate}
    \item $\sF$ is  a linear pullback of a rational foliation either of type $\rat{1,2,2}{n}$ or  $\rat{1,1,3}{n}$;
    \item $\sF \in \logg{1^{n-k+1},2}{k}{n}$ is a logarithmic foliation of type $(1^{n-r +1}, 2)$;
    \item $\sF \in \logg{1^{n+1}}{k}{n}$ is a logarithmic foliation of type $(1^{n+1})$;
    \item  $\sF \in \exc{r_1, \dots,r_m}{k}{n}$ is an  Affine-Lie foliation; or
    \item $\sF \in \lpb{2}{k}{n}$ is a linear pullback of a foliation by curves in $\p{n-k+1}$.
\end{enumerate}
%In particular, the space of foliations of degree $2$ and dimension $k$   has$p(n-k+3) +1$ irreducible components if $n>4$, and $7$irreducible components if $n=4$, where $p(n-k+3) +1$ is the number of partitions of $n-k+3$. 
\end{thmintro}

 




%{\color{red}\begin{thmintro}{} Let $\sF$ be a degree-$2$ foliation on $\pn$ of dimension $ n-2$. Then one of the following holds:\begin{enumerate} \item Then $\sF$ is either a rational foliation of type $\rat{1,2,2}{n}$ or  $\rat{1,1,3}{n}$; \item $\sF \in \logg{1,1,1,2}{k}{n}$ is a logarithmic foliation of type $(1^{n-r +1}, 2)$; \item $\sF \in \logg{1^{n+1}}{n-2}{n}$ is a logarithmic foliation of type $(1^{n+1})$; \item  $\sF \in  \exc{ (4,1^{n-3})}{2}{n}$; $\sF \in \exc{3,2,1^{n-4}}{2}{n}$ or$\sF\in \exc{2^{3} ,1^{n-5}}{2}{n}$\item $\sF \in \lpb{2}{n-2}{n}$ is a linear pullback of a foliation by curves in $\p{3}$.\end{enumerate}In particular, the space of foliations of degree $2$ and codimension $2$ has$8$ irreducible components if $n>4$, and $7$ irreducible components if $n=4$.\end{thmintro}}




The notation is explained in \S\ref{sectionfol}. Theorem \ref{thm:mainA} gives a precise description of the irreducible components of $\Fol(2,k,n)$ whose general element is not algebraically integrable with either $k>2$ and semistable tangent sheaf; or $k=2$ and stable tangent sheaf. Suppose $\sF$ is an algebraically integrable foliation of degree $2$ defined by a rational $(n-k)$-form without zeros in codimension one. In that case, it is known that $\sF$ belongs to one of the rational components: $\rat{1^{n-k-1},2,2}{n}$ or $\rat{1^{n-k},3}{n}$, see Remark \ref{rem:Integraldeg2}. To decide whether every foliation in item (1) of Theorem \ref{thm:mainA}, with $k\leq n-2$, belongs to one of the rational components is a challenging problem whose solution still eludes us. For $k=n-1$ it is true, see \cite[Proposition 3.5]{LPT2}. 

The  Affine-Lie foliations in item (4) of the theorem come from actions of the Lie algebra of affine transformations of the line $\mf{aff(\CC)}$, and these correspond to integer partitions $r_1 + \dots + r_m = n-k+3$. We refer to Pereira-dos Santos \cite{PdS} for a comprehensive description of such foliations. The number of irreducible components of $\Fol(2,k,n)$ grows with the number $p(n-k+3)$ of partitions of $n-k+3$:
\[
\#\{ \text{irreducible components of } \Fol(2,k,n)\} \geq p(n-k+3) +1,
\]
see \cite[Corollary D]{PdS}. 

Along our way to proving Theorem \ref{thm:mainA} we proved the following result, see Proposition \ref{prop:nonsplit}. We observe that this result was also obtained by da Costa \cite{C} with independent techniques.

\begin{propintro}{B}\label{prop:B}
Let $\sF$ be a $2$-dimensional foliation of degree $d$ on $\pn$ such that $h^0(\tF)>0$. Then, either $\tF\cong \opn(a)\oplus \opn(1+a-d)$, $a\in\{0,1\}$, or $\tF$ is not locally free. 
\end{propintro}


Related to holomorphic foliations is the study of Poisson structures, see \S\ref{sectionPoisson} for a definition. A Poisson structure of generic rank $2$ on $\pn$ corresponds to a foliation of dimension $2$ and degree at most $2$. In this framework, Polishchuk found in \cite{Polishchuk} a partial classification of foliations of degree $2$ on $\p3$. Coming from the foliations side, Loray, Pereira, and Touzet \cite{LPT2} classified Poisson structures on Fano $3$-folds of Picard rank one. In \cite[section 5.3]{Matviichuk-Pym-Schedler}, Matviichuk, Pym, and Schedler classified Poisson structures on even-dimensional projective spaces which are log symplectic along the divisor $H$ given by the union of the coordinate hyperplanes, i.e., generically symplectic Poisson structures whose the degeneracy locus is the Poisson divisor $H$.

 

Let $(\pn, \sigma)$ be a holomorphic Poisson structure such that $\sigma \wedge \sigma =0$. 
Suppose that the associated symplectic foliation $\sF_\sigma$  has no divisorial zeros. It follows from Theorem \ref{thm:mainA} that the tangent sheaf $T\sF_\sigma$ of the associated symplectic foliation is either split, semi-stable and logarithmic, or stable with algebraic leaves. This means that either the bivector $\sigma$ is globally decomposable or non-decomposable with some  Poisson divisor $D$. 
In general, we will say that a non-decomposable bivector generic rank 2  is of the \textit{normal crossing type}, if so is $D$.








Due to our study of foliations, we prove the following result about the geometry of Poisson structures with generic rank 2. 




\begin{corintro}{C}\label{Cor:Poisson}
Let $(\PP^{4}, \sigma)$ be a general holomorphic Poisson structure such that
$\sigma \wedge \sigma = 0$. Suppose that the associated symplectic foliation
$\sF_\sigma$ has no divisorial zeros.

If $\sigma$ is decomposable, then either
\begin{itemize}
  \item $\sigma = v_0 \wedge v_1 \in \wedge^2 \mathfrak{h}$, where
  $\mathfrak{h} = H^0(\PP^{4}, T\PP^{4}(-\log D))$ is the abelian Lie algebra
  of vector fields tangent to an arrangement of Poisson planes
  $D = \bigcup_{i=1}^{5} H_i$, and the vector fields $v_i$ induce on each
  symplectic leaf of $\sigma$ a transversely Euclidean holomorphic foliation.

  \item $\sigma = v_0 \wedge v$, where $v_0$ is a constant vector field,
  $v$ is a homogeneous polynomial vector field of degree $2$, $[v_0,v]=0$,
  and $H$ is the Poisson hyperplane determined by
  $v_0 \in \PP H^0(\PP^{4}, T\PP^{4}(-1)) \cong \check{\PP}^{4}$.

  \item $\sigma = v_0 \wedge v_1 \in \wedge^2 \mathfrak{h}$, where
  $\mathfrak{h} = H^0(\PP^{4}, T\PP^{4})$ and $v_0$ and $v_1$ are linear
  vector fields such that $[v_0,v_1]=v_0$, and the vector field $v_0$ induces
  on each symplectic leaf of $\sigma$ a transversely affine holomorphic foliation.
  Moreover, up to a linear change of homogeneous coordinates
  $[x_0:x_1:x_2:x_3:x_4]$ on $\PP^4$, the pair $(v_0,v_1)$ is of one of the
  following two types:

  \begin{itemize}
    \item[\textup{(I)}] \emph{Type $(4,1)$.}
    \[
      v_0
      =
      x_0\,\frac{\partial}{\partial x_1}
      +
      x_1\,\frac{\partial}{\partial x_2}
      +
      x_2\,\frac{\partial}{\partial x_3},
    \]
    \[
      v_1
      =
      x_1\,\frac{\partial}{\partial x_1}
      +
      2x_2\,\frac{\partial}{\partial x_2}
      +
      3x_3\,\frac{\partial}{\partial x_3},
    \]
    so that $[v_0,v_1]=v_0$ and the Jordan type of $v_0$ is $(4,1)$.

    \item[\textup{(II)}] \emph{Type $(3,2)$.}
    \[
      v_0
      =
      x_0\,\frac{\partial}{\partial x_1}
      +
      x_1\,\frac{\partial}{\partial x_2}
      +
      x_3\,\frac{\partial}{\partial x_4},
    \]
    \[
      v_1
      =
      x_1\,\frac{\partial}{\partial x_1}
      +
      2x_2\,\frac{\partial}{\partial x_2}
      +
      x_4\,\frac{\partial}{\partial x_4},
    \]
    so that $[v_0,v_1]=v_0$ and the Jordan type of $v_0$ is $(3,2)$.
  \end{itemize}
\end{itemize}

If $\sigma$ is not decomposable and of normal crossing type, then either
\begin{itemize}
  \item $\sigma \in H^0\bigl(\PP^{4}, \wedge^2 T_{\PP^{4}/\PP^{2}}(-\log D)\bigr)$,
  where $D = H_1 \cup H_2 \cup \mathcal{V}$ is a Poisson divisor such that
  $H_1$ and $H_2$ are hyperplanes and $\mathcal{V}$ is a cubic, or

  \item $\sigma \in H^0\bigl(\PP^{4}, \wedge^2 T_{\PP^{4}/\PP^{2}}(-\log D)\bigr)$,
  where $D = H_1 \cup \mathcal{Q}_1 \cup \mathcal{Q}_2$ is a Poisson divisor
  such that $H_1$ is a hyperplane and $\mathcal{Q}_1, \mathcal{Q}_2$ are quadrics, or

  \item $\sigma \in H^0\bigl(\PP^{4}, \wedge^2 T\PP^{4}(-\log D)\bigr)$,
  where $D = H_1 \cup H_2 \cup H_3 \cup \mathcal{Q}$ is a Poisson divisor such that
  $H_i$ are hyperplanes and $\mathcal{Q}$ is a quadric, and there exists
  $v \in H^0(\PP^{4}, T\PP^{4}(-\log D))$ which is Hamiltonian with respect to $\sigma$.
\end{itemize}
\end{corintro}







 
Consider the space of Poisson structures of generic rank 2 
$$
\mathrm{Poisson}(X)_2 := \left\{ \sigma \in \mathbb{P}H^0(X, \wedge^2 TX) \, \Big| \, [\sigma,\sigma]=0, \sigma \wedge \sigma=0 \, \right\} \subset \mathrm{Poisson}(X).
$$
As a consequence of the study of foliations of degree 2 and dimension 2 we obtain. 
\begin{corintro}{D}
The space $\mathrm{Poisson}(\mathbb{P}^4)_2$ has at least seven irreducible components.
\end{corintro}



This work is organized as follows. In Section \ref{sectionfol} we establish notation and recall basic results on holomorphic foliations. We describe in \S\S\ref{sec:lie-alg-fols} the $2$-dimensional foliations associated with an action of $\mf{aff(\CC)}$, see Theorem \ref{thm:lie-comps}. We call them Affine-Lie foliations. We also describe a rational first integral for a general foliation associated with the partition $(n+1)$, see Proposition \ref{prop:aff-alg-int}. In Section \ref{sec:subfol}, we give a detailed study of subfoliations, especially those arising by strict semistability. In particular, we prove Proposition \ref{prop:B} (see Proposition \ref{prop:nonsplit}) and prepare for the proof of Theorem \ref{thm:mainA}. In Section \ref{sec:main}, we prove Theorem \ref{thm:mainA} by proving two results:  Theorem \ref{thm:dim2-structure} for 2-dimensional foliations, and Theorem \ref{thm:dim>2} for higher dimensions. Finally, in Section \ref{sectionPoisson}, we discuss Poisson structures and we prove Corollary \ref{Cor:Poisson}. 




\subsection*{Acknowledgments}
We thank Jorge Vit\'orio Pereira for his interest in this work and for numerous insightful conversations. MC is partially supported by the Universit\`a degli Studi di Bari and by the PRIN 2022MWPMAB- "Interactions between Geometric Structures and Function Theories" and is a member of INdAM-GNSAGA. AM was partially supported by INCTmat/MCT/Brazil, CNPq grant number 160934/2022-2. 

 
 
 
 

 

 
 
\section{Holomorphic foliations}\label{sectionfol}
In this section, we review the basic concepts needed throughout the work. Let $X$ be a smooth complex projective variety;  an experienced reader may find no difficulty relaxing this requirement. A \emph{singular holomorphic distribution} $\sF$ on $X$ is the data of a saturated subsheaf $\tF$ of the {tangent sheaf} $\tX$ of $X$. These sheaves give a short exact sequence
\begin{equation}\label{eq:dist-def}
    0  \lra  \tF  \lra  \tX  \lra  \nF  \lra  0
\end{equation}
where $\nF$ is a torsion-free sheaf, called the \emph{normal sheaf} of $\sF$, and $\tF$ is a reflexive sheaf, called the \emph{tangent sheaf} of $\sF$. The dimension of $\sF$ is the generic rank of $\tF$. If $\tF$ is involutive, i.e. $[\tF, \tF] \subset \tF$, then $\sF$ is a foliation; this means that through a general point in $X$ passes a unique connected immersed (analytic) subvariety $L \subset X$ such that ${\rm T}L = \tF|_L$, called a \emph{leaf}. A foliation $\sF$ is \emph{algebraically integrable} if its leaves are embedded algebraic varieties. 

Taking exterior power on \eqref{eq:dist-def} we get a morphism $\omega_{\sF}\colon \extp^p\tX \to \det(\nF)$, where $p = \rk \nF$, whose image is $\det(\nF)$ tensored by the ideal sheaf $\sI_Z$ of the vanishing locus of $\omega_{\sF}$ seen as a global section of $\Omega_X^p\otimes \det(\nF)$. We call $Z$ the \emph{singular scheme} of $\sF$, and denote it by $\sing(\sF)$. Since $\tF$ is saturated, $\codim(\sing(\sF))\geq 2$.

Given a distribution $\sF$, the twisted $p$-form $\omega_\sF\in H^0(\Omega_X^p\otimes \det(\nF))$ is rather special, for any $p \in X \setminus \sing(\sF)$ there exists an open neighborhood $U\ni p$ and $1$-forms $\eta_1, \dots, \eta_p \in \Omega_X^1(U)$ such that 
\[
\omega_\sF|_U = \eta_1\wedge \dots \wedge \eta_p.
\]
Moreover, $\sF$ is a foliation if and only if
\[
d\eta_j \wedge \omega_\sF|_U = 0, \, \text{for} \, j= 1, \dots, p.
\]
%%-------------------------

\subsection{Foliations on projective spaces}\label{sect:folsonPn}
A $k$-dimensional distribution $\sF$ on $\pn$ corresponds to a twisted $(n-k)$-form $\omega \in H^0(\Omega_{\pn}^{n-k}(d+n-k+1)$, where $d\geq 0$ is called the \emph{degree} of $\sF$. The Euler sequence allows us to interpret $\omega$ as a polynomial differential form
\[
\omega = \sum_{i=0}^n A_idx_i
\]
where $A_i$ are homogeneous of degree $d+1$, and the contraction with the radial vector files $\rad = x_0\del{ x_0} + \cdots + x_n\del{ x_n}$ vanishes: $\iota_{\rad} \omega = 0$. The local decomposability translates to the Pl\"{u}cker conditions  
\begin{equation}\label{eq:lds}
    ( {i}_v\omega) \wedge \omega = 0 \quad \text{ for every }v
 \in \wedge^{n-k-1}(\mathbb C^{n+1}) ,
\end{equation}
and the integrability condition becomes
\begin{equation}\label{eq:integrability}
     ( {i}_v\omega) \wedge d\omega = 0 \quad \text{ for every }v
 \in \wedge^{n-k-1}(\mathbb C^{n+1}) .
\end{equation}
We refer to \cite{Jou, CPV1} for details. 

Given integers $d$, $k$, and $n$, the conditions \eqref{eq:lds} and \eqref{eq:integrability} define the following subvariety of $\PP H^0(\Omega_{\pn}^{n-k}(n-k+d+1))$:
\begin{equation}\label{eq:def-Fol}
    \Fol(d,k,n) =  \Big\{[\omega]   \mid \omega \text{ is integrable, and } \codim( \sing \omega) \geq 2 \Big\},
\end{equation}
the space of $k$-dimensional foliations of degree $d$ on $\pn$. For $d \leq 1$ the irreducible components of $\Fol(d,k,n)$ are well-known. A degree-$0$ foliation is defined by a linear rational map $\pn \dashrightarrow \p{n-k}$, see \cite[ Th\'eor\`eme 3.8]{CD2}. In particular,
\[
\Fol(0,n-k,n) \cong G(n-k+1, n+1)
\]
the Grassmannian of $(n-k+1)$-planes in $\CC^{n+1}$. For degree one, there is the following classification. 

\begin{theorem}\cite[ Theorem 6.2]{LPT2}\label{thm:fols-deg1}
If $\sF$ is a foliation of degree $1$ and codimension $q$ on $\pn$, then we are in one of the following cases:
\begin{enumerate}
\item $\sF$ is defined by a dominant rational map $\rho \colon \pn \dashrightarrow \PP(1^{(q+1)},2)$, with irreducible general fiber determined by $q$ linear forms and one quadratic form; or
\item $\sF$ is the linear pullback of a foliation induced by a global holomorphic vector field on  $\mathbb{P}^{n-k+1}$. 
\end{enumerate}
\end{theorem}
Hence, $\Fol(1,n-k,n)$ has two irreducible components. In degree $2$, only the case of codimension-one foliations is known. 

\begin{theorem}\cite[Theorem A]{CLN}
The space of codimension-one foliations of degree $2$ in $\pn$, $n\geq 3$, has $6$ irreducible components: 
\begin{itemize}
    \item ${S(2,n)}$ composed of linear pullbacks of degree-$2$ foliations on $\p2$;
    \item ${R(2,2)}$ given by pencils of quadrics;
    \item ${R(1,3)}$ given by pencils of cubics with a triple hyperplane;
    \item ${L(1,1,1,1)}$ whose general member is defined by a logarithmic $1$-form with poles in $4$ hyperplanes;
    \item ${L(1,1,2)}$ whose general member is defined by a logarithmic $1$-form with poles in two hyperplanes and a quadric;
    \item ${E(n)}$ the \emph{exceptional component} whose general member is defined by an algebraic action of the Lie algebra of affine transformations $\mathfrak{aff}(\CC)$.
\end{itemize}
\end{theorem}

Below, we detail each one of these components. We describe their generalizations to any codimension. 

\begin{example}[Linear pullback foliations]\label{ex:lpb-fols}
Given a linear projection $\rho \colon \pn \dashrightarrow \p{n-k+1}$ and a $\sG$ degree-$d$ foliation by curves on $\p{n-k+1}$, then we define $\sF = \rho^*\sG$ a degree-$d$ foliation of dimension $k$. Varying $\rho$ and $\sG$, we get an irreducible component $\lpb{d}{k}{n} \subset \Fol(d,k,n)$. For a proof, see \cite[Corollary 5.1]{CP1}.
\end{example}

\begin{example}[Rational foliations]\label{ex:rat-fols}
Given homogeneous polynomials $F_0, \dots , F_{n-k}$, we have a rational map $\rho \colon \pn \dashrightarrow \PP(d_0, \dots, d_{n-k})$, where $d_j = \deg(F_j)$. Its fibers define an algebraically integrable foliation that we call \emph{rational} if $\rho$ does not ramify in codimension one, i.e., the maximal minors of the Jacobian matrix of $\rho$ do not share a common factor. These foliations fill an irreducible component 
\[
\rat{d_0, \dots, d_{n-k}}{n} \subset \Fol( d, k, n),
\]
where $d = -n+k-1 + \sum_{j=0}^{n-k} d_j$. For a proof, see \cite{CPV1}.    
\end{example}


\begin{remark}\label{rem:Integraldeg2}
In degree $2$, there are two components of rational foliations. Indeed, the equation $2 = -n+k-1 + \sum_{j=0}^{n-k} d_j$ admits only two solutions:
\[
(d_0, \dots , d_{n-k}) = (1^{n-k-2}, 2,2) \text{ or } (1^{n-k-1}, 3).
\]
If a foliation $\sF$ of dimension $k$ and degree 2 is tangent to a rational map $\rho \colon \pn \dashrightarrow \mathbb{P}(d_0, \dots, d_{n-k})$, such that the $\sF$-invariant divisor $\{F_0 = \dots = F_{n-k} = 0\}$ is normal crossing, then by \cite{BruM} we get 
$d_0 + \cdots + d_{n-k} \leq n - k + 3$. This implies that $\sF$ is rational.
\end{remark}

\begin{example}[Logarithmic foliations]\label{ex:log-fols}
A logarithmic foliation on $\pn$ of dimension $k$ is defined by a closed rational $(n-k)$-form
\[
\omega = \sum_{0\leq i_1< \cdots <i_{n-k}\leq r} \lambda_{i_1, \dots , i_{n-k}}\frac{dF_{i_1}}{F_{i_1}}\wedge \cdots \wedge \frac{dF_{i_{n-k}}}{F_{i_{n-k}}}
\]
where $F_0, \dots , F_{r}$ are homogeneous polynomials of degrees $d_0, \dots, d_r$, with $r \geq n-k$, and $\lambda_{i_1, \dots , i_{n-k}} \in \CC$ satisfy Pl\"ucker conditions. We also assume that $\omega$ does not vanish in codimension one. Owing to 
\cite[Theorem 1.17]{CLNlog}, if $r\geq n-k+1$ these foliations form an irreducible component 
\[
\logg{d_0, \dots, d_r}{k}{n}\subset   \Fol( d, k, n), 
\]
where $d = -n+k-1 + \sum_{j=0}^{r} d_j$. We remark that for $r=n-k$ these are rational foliations. 
\end{example}


\begin{remark}\label{rem:log(1,2)}
In degree $2$, the equation $2 = -n+k-1 + \sum_{j=0}^{r} d_j$ has only two solutions with $r\geq n-k+1$: 
\[
(d_0, \dots , d_{r}) = (1^{n-k+1}, 2) \text{ or } (1^{n-k+3}).
\]
Thus $\Fol(2,k,n)$ has two components of logarithmic type. We observe that a foliation of dimension $k$ is a general logarithmic foliation of type $(1^{n-k+1},2)$ if and only if it is the pullback, via a rational map $\rho\colon \pn \dashrightarrow \PP(1^{n-k+1},2)$, of a foliation defined by a global vector field. The proof is a straightforward computation that we leave to the reader. 

A logarithmic foliation of type $(1^{n-k+3})$ and dimension $k>2$ is a linear pullback of a $2$-dimensional logarithmic foliation of the same type. Indeed, for $k>2$, $n-k+3 < n+1$, and, up to a linear change of coordinates, the foliation is defined by fewer than $n+1$ variables. For dimension $k=2$, a general logarithmic foliation of type $(1^{n+1})$ is defined by a pair commuting global vector fields on $\pn$, i.e., it is associated with an action of the commutative Lie algebra $\CC^2\hookrightarrow H^0(\tpn)$. 
\end{remark}

\subsection{Foliations associated with a Lie algebra action}\label{sec:lie-alg-fols}
Given a $k$-dimensional Lie subalgebra $\mf{g} \subset \mf{sl}_{n+1} =  H^0(\tpn)$, it defines an element $\omega(\mf{g}) \in H^0(\Omega_{\pn}^{n-k}(n+1))$. If $\omega(\mf{g})\neq 0$ does not vanish in codimension one, it defines a foliation of degree equal to its dimension, $d=k$. We refer to \cite{CD2,CP1,PdS} for generalities and details about these foliations. Here we will focus on $2$-dimensional Lie algebras. In \cite{PdS}, Pereira and dos Santos describe the scheme of Lie subalgebras and the irreducible components of the space of foliations they define.


%We reproduce below some of their arguments.


Consider the algebraic variety of $2$-dimensional Lie subalgebras of $\sln_{n+1}$: 
\[
{\rm SLie }(2, \sln_{n+1}) = \{\,V \subset \sln_{n+1} \mid [V,V]\subset V, \, \dim V = 2 \,\}.
\]
It is a closed subvariety of the Grassmannian $G(2, \sln_{n+1})$ defined by the equations imposed by the Lie bracket. ${\rm SLie }(2, \sln_{n+1})$ has several irreducible components, one $\mathcal{C}$ corresponding to commutative Lie algebras and another for each integer partition of $n+1$:
\[
{\rm SLie }(2, \sln_{n+1}) = \mathcal{C} \cup \bigcup_{r_1+\cdots +r_m = n+1}\mathcal{S}_{(r_1, \dots, r_m)}
\]

%A result by Motzkin-Taussky says that the algebraic variety parametrizing pairs of commuting matrices is irreducible, and it is the closure of the set of pairs of (simultaneously) diagonalizable matrices, see \cite[Theorem 5]{MoTa} and \cite{Gu}. This gives one irreducible component $\mathcal{C}$ of ${\rm SLie }(2, \sln_{n+1})$. The other components correspond to non-Abelian subalgebras.

% Up to isomorphism, there is only one $2$-dimensional noncommutative Lie algebra: $\mathfrak{aff}(\CC) = \langle\, e,f\mid [e,f] = e \,\rangle$. Then, computing the derived subalgebra induces a map 
% \[
% \Phi \colon {\rm SLie }(2, \sln_{n+1})\setminus \mathcal{C}  \lra  \mathcal{N}; \quad \mathfrak{g} \longmapsto [\mathfrak{g}, \mathfrak{g}],
% \]
% where $\mathcal{N} \subset \PP(\sln_{n+1})$ is the projectivization of the nullcone, the subvariety of nilpotent elements. The celebrated Jacobson-Morosov Theorem says that any nilpotent element $e$ of a semisimple Lie algebra extends to an $\sln_2$-triple: $\{f,e,h\}$, $[h,e]=2e$, $[h,f] =-2f$, and $[e,f] = h$ \cite[Theorem 3.7.1]{CG-rep}. Then, $\Phi$ is surjective. 

% The variety $\mathcal{N}$ can be stratified by the orbits of nilpotent elements under the action of ${\rm SL}_{n+1}$ by conjugation. Each nilpotent orbit corresponds to a Jordan normal form, thus to a partition of $n+1$. Given a partition $r_1 + \cdots + r_m = n+1$, we denote the corresponding orbit by $\OO_{(r_1, \dots, r_m)}$. Note that $\Phi^{-1}(\OO_{(r_1, \dots, r_m)}) =  {\rm SL}_{n+1}\cdot \Phi^{-1}([v])$, for any $[v]\in \OO_{(r_1, \dots, r_m)}$. Therefore, 
% \[
% \mathcal{S}_{(r_1, \dots, r_m)} := \overline{\Phi^{-1}(\OO_{(r_1, \dots, r_m)})}
% \]
% are irreducible closed subvarieties of ${\rm SLie }(2, \sln_{n+1})$ whose (finite) union cover ${\rm SLie }(2, \sln_{n+1})\setminus \mathcal{C}$. To see that these are irreducible components is enough to prove they have the same dimension: $\dim \sln_{n+1} -1$. Note that for $[v] \in \mathcal{N}$ we have 
% \[
% \Phi^{-1}([v]) = \{\, \langle v,w \rangle \in {\rm SLie }(2, \sln_{n+1})  \mid [v,w] \in [v] \,\} \cong  \frac{\{\, [w] \in \PP(\sln_{n+1}) \mid [v,w] \in [v]\,\}}{[v]}.
% \]
% Fix $w_0$ such that $[v,w_0]=v$. Then, any $w$ as above can be written as $w= \lambda w_0 + w'$ with $\lambda \in \CC^*$ and $[v,w']=0$, i.e., $w' \in C_{\sln_{n+1}}(v)$, the centralizer of $v$. Thus, $\dim \Phi^{-1}([v]) = \dim C_{\sln_{n+1}}(v) -1 $. On the other hand, $\dim \OO_{(r_1, \dots, r_m)} = \dim \sln_{n+1} - \dim \dim C_{\sln_{n+1}}(v)$. Hence, 
% \[
% \dim \mathcal{S}_{(r_1, \dots, r_m)}  = \dim \sln_{n+1} -1.
% \]


Consider the map $\mathcal{F} \colon {\rm SLie }(2, \sln_{n+1})   \lra  \Fol(2,n-2,n)$ associating with a Lie subalgebra $\mathfrak{g}$ the foliation $\mathcal{F}(\mathfrak{g})$ it defines, whenever it makes sense. 
%It follows from \cite[Theorem A]{Ve} that $\mathcal{F}$ is an embedding where it is defined. 
We saw in Remark \ref{rem:log(1,2)} that $\mathcal{F}(\mathcal{C}) = {\rm Log}(1^{n+1};2,n)$ the logarithmic component. For the other components, there is the following result, which is a restatement of \cite[Theorem 3.6]{PdS}. They only consider partitions with one part at least $5$; however, the argument works for other partitions with minor changes. 

\begin{theorem}\cite[Theorem 3.6]{PdS}\label{thm:lie-comps}
We have that  $\exc{r_1, \dots, r_m}{2}{n} :=  \mathcal{F}(\mathcal{S}_{(r_1, \dots, r_m)})$ is an irreducible component of $\Fol(2,n-2,n)$ unless $(r_1, \dots, r_m)$ is one of the following partitions:
\[
\{ (3,1^{n-2}), (2,2,1^{n-3}), (2, 1^{n-1}), (1^{n+1}) \}.
\]
Moreover, $\mathcal{F}$ is not defined on $\mathcal{S}_{(1^{n+1})}$ and $\mathcal{S}_{(2, 1^{n-1})}$.
\end{theorem}



% \begin{proof}
% We may assume $n\geq 4$ since the case $n=3$ is well-known. Then \cite[Corollary 6.1]{CP1} says that a small deformation of a foliation $\mathcal{F}(\mathfrak{g}) \in \mathcal{F}(\mathcal{S}_{(r_1, \dots, r_m)})$ is induced by another Lie algebra $\mathfrak{g}' \in {\rm SLie }(2, \sln_{n+1})$ if $\codim \sing (\mathcal{F}(\mathfrak{g})) \geq 3$. Therefore, it is enough to produce a foliation in $\mathcal{F}(\mathcal{S}_{(r_1, \dots, r_m)})$ with singularities in codimension at least $3$.

% Recall that the singular locus of a distribution given by a pair of vector fields $v = \sum_{j=0}^n a_j \del{ x_j}$ and $w = \sum_{j=0}^n b_j \del{ x_j}$ is defined by the homogeneous ideal $I_3(A)$ generated by the $3\times 3$ minors of the matrix
% \[
% A = \begin{bmatrix}
%     a_0  &  \dots & a_n \\
%     b_0  &  \dots & b_n \\
%     x_0 & \dots & x_n
% \end{bmatrix}.
% \]
% Given any submatrix $A'$ of $A$, we have $I_3(A) \subset I_3(A')$. Hence, to show that $\codim V(I_3(A)) \geq 3$ is enough to prove that the same holds for some submatrix $A'$. We will apply this strategy to our specific case. 

% Fix a partition $r_1 \geq r_2 \geq \dots \geq r_m$ and fix adapted homogeneous coordinates:
% \[
% (x_0^1: \dots: x_{r_1-1}^1: \dots : x_0^k : \dots : x_{r_m-1}^k).
% \]
% Then define the vector fields $v_{(r_1, \dots, r_m)} = \delta_{r_1} + \cdots + \delta_{r_m} $ and $w_{(r_1, \dots, r_m)} = \rho_{r_1} + \cdots + \rho_{r_j}$, where
% \begin{equation}\label{eq:basic-nilp}
%     \delta_{r_j} = \sum_{i=0}^{r_j-2} x_{i}^j \del{ x_{i+1}^j}, \quad \text{and} \quad \rho_{r_j} = \sum_{i =0}^{ r_j-1} (r_j-1 - 2i + \epsilon_j) x_i^j\del{ x_{i}^j}, \quad \epsilon_1+ \dots +  \epsilon_k = 0 \in \CC.
% \end{equation}
% It follows that $[w_{(r_1, \dots, r_m),} v_{(r_1, \dots, r_m)}] = v_{(r_1, \dots, r_m)}$. Let $A$ be the matrix defined above. If either $r_1 \geq 5$, or $r_1 = 4$ and $r_2 \geq 1$, or $r_1 = 3$ and $r_2 \geq 2$, then let $A'$ be the submatrix given by the first $5$ columns of $A$. If $r_1 = r_2 = r_3 = 2$ then let $A'$ be the submatrix composed of the first $6$ columns of $A$. Then, for $\epsilon = (\epsilon_1, \dots, \epsilon_k)$ generic, $\codim V(I_3(A')) \geq 3$. This is a lengthy but straightforward computation, and we leave it to the reader. 

% The partitions not covered by the above argument are: $(1^{n+1})$, $(2,1^{n-1})$, $(3,1^{n-2})$, $(2,2,1^{n-3})$, since $n\geq 4$. Note that $v_{(1^{n+1})}= 0$ and $v_{(2,1^{n-1})} = x_0^1\del{ x_1^1}$ has singularities in codimension one. The other cases have singularities in codimension $2$. 
% \end{proof}

\begin{remark}\label{Alg: log(1,2)} 
The partition $(4)$ gives the classical exceptional component on $\p3$ \cite{CLN}. The partition $(n+1)$ for $n\geq 3$ figures in \cite[Example 6.8]{CP1}. These are all rigid, i.e., the component $\exc{n+1}{2}{n}$ is the orbit closure of a general foliation in it. We also remark that $\mathcal{F}(\mathcal{S}_{(3,1^{n-2})})$  and $\mathcal{F}(\mathcal{S}_{(2,2,1^{n-3})})$ do not fill irreducible components of $\Fol(2,n-2,n)$, both fall into the logarithmic component $\logg{1^{n-1},2}{2}{n}$, see Proposition \ref{prop:class-om-nilp} below.
\end{remark}

\begin{remark}\label{rem:exceptional}
Due to \cite[Corollary 5.1]{CP1}, if $k\geq 3$ and $r_1 + \dots + r_m = n-k+3$, there are irreducible components $\exc{r_1, \dots, r_m}{k}{n} \subset \Fol(2,k,n)$ whose general element is a linear pullback of a foliation in $\exc{r_1, \dots, r_m}{2}{n-k+2}$.  We call these the \emph{Affine-Lie foliations}.   
\end{remark}


\subsubsection{First integrals for  Affine-Lie foliations}\label{normal-form-poisson}
Fix a partition $r_1 \geq r_2 \geq \dots \geq r_m$ and fix adapted homogeneous coordinates:
\[
(x_0^1: \dots: x_{r_1-1}^1: \dots : x_0^k : \dots : x_{r_m-1}^k).
\]
Then define the vector fields $v_{(r_1, \dots, r_m)} = \delta_{r_1} + \cdots + \delta_{r_m} $ and $w_{(r_1, \dots, r_m)} = \rho_{r_1} + \cdots + \rho_{r_j}$, where
\begin{equation}\label{eq:basic-nilp}
    \delta_{r_j} = \sum_{i=0}^{r_j-2} x_{i}^j \del{ x_{i+1}^j}, \quad \text{and} \quad \rho_{r_j} = \sum_{i =0}^{ r_j-1} (r_j-1 - 2i + \epsilon_j) x_i^j\del{ x_{i}^j}, \quad \epsilon_1+ \dots +  \epsilon_k = 0 \in \CC.
\end{equation}
It follows that $[w_{(r_1, \dots, r_m)}, v_{(r_1, \dots, r_m)}] = v_{(r_1, \dots, r_m)}$. Then $v_{(r_1, \dots, r_m)}$ and $w_{(r_1, \dots, r_m)}$ define a foliation in $\exc{r_1, \dots, r_m}{2}{n}$. We will focus on the partition $(n+1)$.


The vector field $\delta_{m}$, as defined in \eqref{eq:basic-nilp}, is sometimes called the basic Weitzenb\"ock derivation. In general, a Weitzenb\"ock derivation is a linear nilpotent derivation. They play an important role in Invariant Theory, and the name comes from Weitzenb\"ock's Theorem (see \cite{Se}) which asserts that any linear action of the additive group $\mathbb{G}_a$ on a polynomial ring has a finitely generated ring of invariants. That is, the ring 
\[
\ker \delta = \{\, f\in k[y_1, \dots,y_n] \mid \delta(f) = 0\,\}
\]
is finitely generated (over $k$) for $\delta$ a Weitzenb\"ock derivation. This fails for polynomial locally nilpotent derivations of higher degree, see \cite[Example 6.3.1]{Now}. In \cite{JS}, the authors give explicit generators of the field of fractions of the constant ring of the basic Weitzenb\"ock derivation. Using this description, we can show that a foliation in $\exc{n+1}{2}{n}$ is algebraically integrable.

\begin{proposition}\label{prop:aff-alg-int}
A general Affine-Lie foliation in $\exc{n+1}{2}{n}$ is algebraically integrable.
\end{proposition}

\begin{proof}
It is enough to prove that the foliation $\sF$ defined by $\delta_{n+1}$ and $\rho_{n+1}$ in \eqref{eq:basic-nilp} is algebraically integrable. Note that we only need to prove that there exist $n-2$ algebraically independent functions $f_i/g_i \in \CC(x_0,\dots, x_n)$ such that $f_i$ and $g_i$ are homogeneous of the same degree, and $v(f_i/g_i)=w(f_i/g_i)=0$.

According to \cite[\S5]{JS}, the field of fractions of $\ker v$ is generated by algebraically independent elements $z_1, \dots, z_n$ given by: 
\begin{align*}
    z_1 & = x_0, \\
    z_{2i} & = \sum_{j=0}^{2i} (-1)^j x_jx_{2i-j}, \, \text{and} \\
    z_{2i+1} & = \frac{x_0}{2n-4i}\left(\sum_{j=0}^{2i}(-1)^j (j+1)(n-j)x_{j+1}x_{2i-j} \right. \\ &  +\sum_{j=0}^{2i}(-1)^j (2i+1-j)(n-2i+j)x_jx_{2i+1-j}\Bigg) - x_1\sum_{j=0}^{2i}(-1)^jx_jx_{2i-j}, \, \text{for} \, j\geq 1.
\end{align*}
One may verify that $w(z_1) = \frac{n}{2}z_1$, $w(z_{2i}) = n-2i$, and $w(z_{2i+1}) = \frac{3n}{2}-2i-1$ for $j\geq 1$. A first integral for $\sF$ is then a rational function $f = z_1^{a_1}\cdots z_n^{a_n}$ such that $w(f) = 0$, hence, the exponents $a_j$ satisfy
\[
\begin{cases}
    \displaystyle a_1 + \sum_{i = 1}^{\lfloor n/2 \rfloor} (2a_{2i}+3a_{2i+1}) = 0 \\
    \displaystyle \frac{n}{2}a_1 + \sum_{i = 1}^{\lfloor n/2 \rfloor} \left( (n-2i)a_{2j}+ \left(\frac{3n}{2}-2i-1 \right)a_{2j+1}\right) = 0
\end{cases}.
\]
The $2\times n$ matrix $A$ associated with this system is of rank $2$. Then we can find $n-2$ integer-valued solutions, linearly independent over $\mathbb{Q}$, which provide $n-2$ algebraically independent rational first integrals for $\sF$. Explicitly, an independent set of first integrals is:
\[
\left\{ \frac{z_{2+i}^2z_1^{4\lfloor i/2 \rfloor}}{z_2^{i+2}}, \, i = 1, \dots , n-2  \right\}.
\]
for $i$ even one can take a square root.
\end{proof}

From the above proof, we also conclude that $\sF$ is given by the $(n-2)$-form
\[
\omega = \frac{1}{z_1^{\frac{n^2- n-4}{2}}}\iota_{\rad} dz_2 \wedge dz_3 \wedge d(z_1^{2}z_{4}) \wedge d(z_1^{2}z_{5}) \wedge d(z_1^{4}z_{6}) \wedge \dots \wedge d(z_1^{2\lfloor \frac{n-2}{2} \rfloor}z_{n}).
\]

\section{Study of subfoliations}\label{sec:subfol}
In this section, we study subfoliations $\sG \subset \sF$ of a fixed foliation $\sF$, with special attention to subfoliations defined by actions of Lie algebras. A Lie algebra naturally attached to a foliation $\sF$ on $\pn$ is given by its infinitesimal automorphisms
\[
\aut(\sF) = \{\, v \in H^0(\tpn) \mid \exists c\in \CC \text{ s.t. } \lieder_v\omega =c\,\omega \,\},
\]
where $\lieder_v$ stands for the Lie derivative in the direction of $v$ and $\omega$ is a twisted differential form defining $\sF$. $\aut(\sF)$ is the Lie algebra of the closed algebraic subgroup $\Aut(\sF) \subset \PGL(n+1,\CC)$ that preserves $\sF$. The global vector fields tangent to $\sF$ belong to $H^0(\tF)$ and define a Lie subalgebra of $\aut(\sF)$, usually denoted by $\fix(\sF)$. The exponential of $\fix(\sF)$ is not necessarily a closed subgroup. Let $\afix(\sF)$ denote the Lie algebra of the smallest algebraic subgroup whose Lie algebra contains $\fix(\sF)$. Note that
\[
\fix(\sF) \subset \afix(\sF) \subset \aut(\sF).
\]
The computation of $\afix(\sF)$ may be carried out as in \cite[Chapter II, \S7]{Borel}. 

 
Another way to obtain special subdistributions/subfoliations is slope-stability. Given a coherent sheaf $F$ on $\pn$ of positive rank $\rk F$, the slope is defined by $\mu(F) :=c_1(F)/\rk F$. Then $F$ is called semistable if $\mu(F) \geq \mu(E)$ for every $E\subset F$ such that $\rk E <\rk F$. $F$ is called stable if the inequality is always strict. One also defines
\begin{align*}
    \mu_{\rm max}(F) &= \max \{\, \mu(E) \mid E \subset F \,\}; \\
    \mu_{\rm min}(F) &= \min \{\, \mu(Q) \mid F \twoheadrightarrow Q \, \text{and } Q \text{ is torsion free}  \,\}.
\end{align*}

We start with the following general observation that we will apply to $\fix(\sF)$. Note that the condition $\mu_{\rm max}(F) \leq 0$ holds for $F$ semistable with $c_1(F) = 0$. 

\begin{lemma}\label{lem:sstab-h0}
Let $F$ be a rank-$r$ torsion-free sheaf on $\pn$ such that $\mu_{\rm max}(F) \leq 0$. Then the evaluation morphism is injective
\[
H^0(F) \otimes \opn \hookrightarrow F
\]
In particular, $h^0(F) \leq r$. If equality holds, then $F = \opn^{\oplus r}$. 
\end{lemma}

\begin{proof}
Let $h = h^0(F)$ and consider the evaluation map $\phi \colon \opn^{\oplus h} \to F$ generated by linear independent global sections. We get a short exact sequence
\[
0 \lra \ker \phi \lra \opn^{\oplus h} \lra \im \phi \lra 0.
\]
Note that $\im \phi$ must be torsion free. Hence, $\ker \phi \subset \opn^{\oplus h}$ is saturated. Since $\opn^{\oplus h}$ is semistable, $\mu(\ker \phi) \leq 0$ and $\mu(\im \phi) \geq 0$. On the other hand, $\mu(\im \phi) \leq \mu_{\rm max}(F) \leq 0$. Then $\mu(\im \phi) = \mu(\ker \phi) = 0$. By \cite[Corollary 1.6.11]{HL-sheaves}, $\ker \phi$ must be a direct summand of $\opn^{\oplus h}$. One can interpret the map $\ker \phi = \opn^{\oplus a} \hookrightarrow \opn^{\oplus h}$ as $\CC$-linear relations among the global sections defining $\phi$. Therefore, $\ker \phi = 0$, i.e., $\phi$ is injective, and $h\leq \rk F$. 

Suppose $h = \rk F$ and consider the exact sequence
\[
0 \lra \opn^{\oplus \rk F} \lra F \lra \coker\phi \lra 0.
\]
Then $\coker\phi$ is a torsion sheaf whose first Chern class is a positive multiple of the divisorial part of its support. On the other hand, $\mu(F) \leq \mu_{\rm max}(F) \leq 0$. Then $c_1(Q) = c_1(F)  \leq 0$, which implies that $Q$ is supported in codimension at least 2. Since $F$ is reflexive, $\phi$ must be an isomorphism. 
\end{proof}


\begin{proposition}\label{prop:sstable-subfol}
Let $\sF$ be a foliation on $\pn$ such that $\mu_{\rm max}(\tF) \leq 0$, then it admits a subfoliation of dimension $h^0(\tF)$ generated by global vector fields.
\end{proposition}

\begin{proof}
By Lemma \ref{lem:sstab-h0}, the evaluation map  $\phi \colon H^0(\tF) \otimes \opn \to \tF$ is injective. Then $H^0(\tF)\otimes \opn \hookrightarrow \tF \hookrightarrow \tpn$ defines a subdistribution, which is a foliation since $H^0(\tF) = \fix(\sF)$ is closed under the Lie bracket. 
\end{proof}





Next, we provide a simple technical observation that will be useful to constrain the singular locus of a subfoliation $\sG$ whose tangent sheaf is locally free. Exact sequences, such as the one below, are known in Commutative Algebra as Bourbaki sequences; see, for instance, \cite{HKS} and the references therein. 

\begin{lemma}\label{lem:codim-ext}
Consider an exact sequence of sheaves on $\pn$ (or any smooth variety)
\[
0 \lra E \lra F \lra \sI_W(a) \lra 0
\]
where $E$ is locally free, $F$ is reflexive, and $\codim W \geq 2$. Then either 
\begin{itemize}
    \item $W= \emptyset$ (hence $F$ is locally free), or
    \item $W$ is pure of codimension two and Cohen-Macaulay off $\sing(F)$. 
\end{itemize}


\end{lemma}

\begin{proof}
Since $F$ is reflexive, it is locally free off $Z = \sing(F)$ whose codimension is at least three. Note that $Z\subset W$ (set-theoretically) since $F$ must be locally free at any point off $W$. At any point $p\in W\setminus Z$, the Hilbert-Burch theorem says that $W$ is Cohen-Macaulay of codimension two. 

We only need to prove that $W$ is pure or, equivalently, that $\OO_W$ satisfies the Serre condition $S_{1,2}$, see \cite[Proposition 1.1.10]{HL-sheaves}. This condition means that $\depth(\OO_{W,p}) \geq 1$ for any point $p\in W$ of codimension at least $3$. Fix such a point $p$. Since $F$ is reflexive and $E$ is locally free, we have $\depth(F_p) \geq 2$ and $\depth(E_p) =\codim p \geq 3$. Then $\depth(\sI_{W,p}) \geq 2$ and $\depth(\OO_{W,p}) \geq 1$, see \cite[Proposition 1.2.9]{BH-CMrings}.
\end{proof}



\begin{proposition}\label{prop:sing2xsplit}
Let $\sF$ be a distribution on $\pn$ such that there exists a subdistribution $\sG \subset \sF$ satisfying: $\dim \sG = \dim \sF - 1$, $\tG$ is locally free and $H^1(\tG(c_1(\tG) - c_1(\tF))) = 0$. Then either 
\begin{itemize}
    \item $\tF \cong \tG \oplus \opn(c_1(\tF) - c_1(\tG))$, or
    \item $\tF/\tG = \sI_W(c_1(\tF) - c_1(\tG))$ with $W \subset \sing(\sG)$ a subscheme of pure codimension two.
\end{itemize} 

% $\tF \cong \tG \oplus \opn(c_1(\tF) - c_1(\tG))$ or $\sG$ has singularities in codimension two.
\end{proposition}


\begin{proof}
By hypothesis, $\tF$ and $\tG$ fit in the following commutative diagram with exact rows and columns, where $\delta = c_1(\tF) - c_1(\tG)$.
\begin{equation}\label{eq:diagsubfol}
\begin{tikzcd}
%& 0 \arrow[d,] & 0 \arrow[d] & & \\
 \tG \arrow[r, equal ] \arrow[d, hook, "\rho"] & \tG \arrow[d, hook, "\phi"] & \\
 \tF \arrow[r,hook, "\psi"] \arrow[d,two heads] & \tpn \arrow[d,two heads] \arrow[r] & \nF \arrow[d,equal] \\
 \sI_W(\delta) \arrow[r,hook]  & \nG \arrow[r]  & \nF 
\end{tikzcd}
\end{equation}
By Lemma \ref{lem:codim-ext}, either $W= \emptyset$ or $W$ is pure of codimension two. If $W = \emptyset$, then $\tF$ corresponds to an extension class in $\Ext^1(\opn(\delta) , \tG) = H^1(\tG(-\delta))$ which is zero by hypothesis. Thus $\tF = \tG \oplus \opn(\delta)$ in this case. 

Assume $W \neq \emptyset$. By Lemma \ref{lem:codim-ext}, $W$ is pure of codimension two. Let $p$ be an associated point of $\OO_W$. We have that $p$ has codimension two and, on the other hand, $\tF$ is reflexive. Thus $\tF_p$ is a free $\OO_{\pn,p}$-module. Of course, the same holds for $\tG$ and $\tpn$. Then, we may regard $\phi_p = \psi_p\rho_p$ as a product of matrices. 

The ideal sheaf $\sI_Z$ of $Z= \sing (\sG)$ is defined by the $s\times s$ minors of $\phi$, where $s = \dim \sG$. Also note that $\sI_{W,p}$ is defined by the $s\times s$ minors of $\rho_p$. Then, by Binet-Cauchy Theorem, \cite[p.11]{Gantm.vol1}, the $s\times s$ minors of a product of matrices is a combination of $s\times s$ minors of both matrices. Hence, $\sI_{Z,p} \subset \sI_{W,p}$. Since $p$ is arbitrary, we get $\sI_Z \subset \sI_W$, see \cite[\S6]{Matsumura}. Therefore, $W\subset \sing \sG$.
\end{proof}

\begin{remark}\label{rem:integ-subfol}
1) Consider $\sF$ and $\sG$ as in the statement of Proposition \ref{prop:sing2xsplit}. Note that, in the first case, $\tF = \tG \oplus \opn(\delta)$ is a subsheaf of $\tpn$, hence $\delta \leq 1$, i.e., $c_1(\tF) \leq c_1(\tG)+1$. This is a simple constraint for having a subdistribution that splits. 

\noindent 2) If $\sF$ is integrable, then the integrability of $\sG$ is measured by the map $\extp^2\tG \to \tF/\tG$ that sends $v\wedge w \mapsto [v,w] + \tG$. Since $\tG$ is locally free, this produces a section in $H^0(\extp^2\tG^\vee(\delta)\otimes \sI_W)$, with $W$ possibly empty. Then, $\sG$ is automatically integrable if this cohomology vanishes.  
\end{remark}


% \begin{remark}
% The last part of the proof of Proposition \ref{prop:sing2xsplit} holds in more generality. If $W$ is as in diagram \eqref{eq:diagsubfol} and has codimension $c$, then let $W_c$ be its pure part of codimension $c$. Then $\inext^j(\sI_W(\delta), \opn) = 0$ for $j < c-1$ and $\inext^{c-1}(\sI_W(\delta), \opn) = \omega_{W_c}(n+1-\delta)$. Thus $|W_c|$ is contained in the support of $\inext^{c-1}(\nG, \opn)$, hence in $\sing(\sG)$. One natural question is whether the scheme $W_c$ is contained in $\sing(\sG)$. 
% \end{remark}


 
Combining Proposition \ref{prop:sstable-subfol} and Proposition \ref{prop:sing2xsplit}, one gets the following result.
\begin{corollary}\label{cor:sub-sing-cod2}
Let $\sF$ be a degree-$d$ foliation of dimension $k$ on $\pn$. Suppose that $h^0(\tF) = k -1$ and $\mu_{\rm max}(\tF) \leq 0$, then either $\tF \cong \opn^{\oplus k-1} \oplus \opn(k-d)$ or $\fix(\sF)$ defines a subfoliation of dimension $k-1$ and with singularities in codimension two.
\end{corollary}


In the next two subsections, we investigate foliations of dimensions 2 and 3 such that $0 < h^0(\tF) \leq \dim \sF -1$.




\subsection{2-dimensional foliations tangent to a global vector field}
For 2-dimensional foliations, $\tF$ is a rank-2 sheaf. Then, Corollary \ref{cor:sub-sing-cod2} applies to any foliation of degree $d\geq 2$ such that  $h^0(\tF(-1)) = 0$. 


\begin{remark}\label{rem:raphael}
A result by da Costa, see \cite[Theorem C]{C}, says that foliations on $\p3$ with locally-free tangent sheaf having a global section must split as a direct sum of line bundles. Using Corollary \ref{cor:sub-sing-cod2} we can extend this result to $2$-dimensional foliations on $\pn$, $n\geq 3$. 
\end{remark}


% {\color{red} Consider $L\hookrightarrow \tF \hookrightarrow TX$ with $X$ smooth, $\dim \sF = k \geq 2$ and suppose $\tF$ is locally free in codimension $k$. Then $\tF^\vee \twoheadrightarrow L^\vee\otimes \sI_Z$ with $\codim(Z) \leq \rk \tF = k$. It follows that $Z$ is contained in the singular locus of the foliation $\sG$ defined by $L\hookrightarrow TX$.

% If $\sG$ has sing in codimension $\geq k+1$ then $Z=\emptyset$ and 
% \[
% 0 \lra Q^\vee \lra \tF^\vee \lra L^\vee \lra 0
% \]
% This gives an alternative to Raphael's Thm 5.4. Not sure it leads to some generalization.
% }

\begin{proposition}\label{prop:nonsplit}
Let $\sF$ be a $2$-dimensional foliation of degree $d$ on $\pn$ such that $h^0(\tF)>0$. Then, either $\tF\cong \opn(a)\oplus \opn(1+a-d)$, $a\in\{0,1\}$, or $\tF$ is not locally free. 
\end{proposition}

\begin{proof}
If $h^0(\tF(-1))>0$ then $\sF$ is a linear pullback and $\tF \cong \opn(1)\oplus \opn(1-d)$. Suppose that $h^0(\tF(-1)) = 0$ and $\tF \not\cong \opn \oplus \opn(2-d)$. Since $\rk \tF = 2$, it is enough to prove that $c_3(\tF) \neq 0$. Consider the exact sequence from the first column of \eqref{eq:diagsubfol}:
\[
0  \lra  \opn  \lra  \tF  \lra  \sI_W(2-d)  \lra  0
\]
By Corollary \ref{cor:sub-sing-cod2} $W$ has pure codimension $2$. Furthermore, if $n>3$, then $W\cong \p{n-2}$ is a linear subspace; and if $n=3$, then $W$ is either one line or a (possibly degenerate) disjoint union of two lines, see \cite[\S2.3]{GJM} and Lemma \ref{lem:veccd2} below. 

If $n\geq 4$, we take a general $L \cong \p3$ linearly embedded in $\pn$ transverse to $W$ then the above sequence restricts to $L$ as
\[
0  \lra  \OO_L  \lra  \tF|_L  \lra  \sI_{W\cap L}(2-d)  \lra  0,
\]
$W\cap L$ is a line, and  $c_i(\tF|_L) = c_i(\tF)$ for $i = 1,2,3$. Therefore, we may assume that $n=3$. 

By \cite[Theorem 4.1]{H2}, $c_1(\tF) = 2-d$, $c_2(\tF) = \deg( W)$ and $c_3(\tF) = 2p_a(W)-2 + \deg(W)(d+2)$. If $W$ is a line then $\deg(W) = 1$ and $p_a(W) = 0$, hence $c_3(\tF) = d \geq 1$. If $W$ is (possibly the limit of) a disjoint union of two lines, $\deg(W) = 2$ and $p_a(W) = -1$, hence $c_3(\tF) = 2d\geq 2$. 
\end{proof}

We classify the possible subfoliations using Jordan normal forms.

\begin{lemma}\label{lem:veccd2}
Let $v\in H^0(\tpn)$ such that $v$ has zeros in codimension two. Then, up to a linear change of coordinates, either
\begin{enumerate}
 \item $v = x_0\del{ x_0} + \lambda x_1\del{ x_1}$ with $\lambda \in \mathbb{C}^*$; 
 \item $v = (x_0+ x_1)\del{ x_0} + x_1\del{ x_1}$;
 \item $v = x_1\del{ x_0} + x_3\del{ x_2}$;
 \item $v = x_1\del{ x_0} + x_2\del{ x_1}$;
 \item { $v = x_1\del{ x_0} + x_2\del{ x_2}$.} 
\end{enumerate}
\end{lemma}

And this allows us to conclude that these subfoliations are algebraically integrable. 

\begin{proposition}\label{prop:sub-deg2-alg}
Let $\sF$ be a foliation of dimension $2$ and degree $d\geq 1$. Suppose that $h^0(\tF(-1)) = 0$ and $h^0(\tF)>0$ and $\tF  \not\cong \opn\oplus \opn(2-d)$. For any $ s \in H^0( \tF)\setminus \{0 \}$, the induced foliation by curves is algebraically integrable with singularities in codimension two.
\end{proposition}



\begin{proof}
Let $\sG$ be the foliation induced by a nonzero section $s \in H^0( \tF)$. By Corollary \ref{cor:sub-sing-cod2}, we know that $\sG$ has singularities in codimension two. Suppose, aiming at a contradiction, that the general leaf of $\sG$ is not algebraic. Then, due to Lemma \ref{lem:veccd2}, $\sG$ is defined by either $v = x_0\del{ x_0} + \lambda x_1\del{ x_1}$ with $\lambda \not \in \mathbb{Q}$, or $v = (x_0+ x_1)\del{ x_0} + x_1\del{ x_1}$, or { $v = x_1\del{ x_0} + x_2\del{ x_2}$}. Let $\omega \in H^0(\omn{n-2}(d+n-1))$ be a $(n-2)$-form defining $\sF$. 

\medskip\noindent
\textbf{Case 1:} Suppose $v = x_0\del{ x_0} + \lambda x_1\del{ x_1}$ with $\lambda \not \in \mathbb{Q}$. Then $x_0 \del{ x_0}, x_1 \del{ x_1} \in \afix(\sF)$:
\[
\lieder_{x_0 \del{ x_0}} \omega = \lambda_0 \omega \quad  \text{and} \quad \lieder_{x_1 \del{ x_1}} \omega = \lambda_1 \omega.
\]
Then $\omega$ is homogeneous of degree $\lambda_j$ on the variable $x_j$, $j =0,1$. Since $H^0(\tF(-1)) = 0$, we have $\lambda_0= \lambda_1 = 1$. Using that $\iota_v \omega =0$, one concludes that 
\[
\omega = (x_0dx_1 - \lambda x_1dx_0)\wedge \alpha + x_0x_1\beta + dx_0\wedge dx_1\wedge \gamma
\]
where $\alpha$, $\beta$, and $\gamma$ do not depend on $x_0,x_1,dx_0,dx_1$. From $\iota_{\rad} \omega = 0$, one gets $\gamma = 0$ and $\alpha = \frac{1}{\lambda -1} \iota_{\rad} \beta$. Thus, up to a constant,
\[
\omega = (x_0dx_1 - \lambda x_1dx_0)\wedge \iota_{\rad}\beta + (\lambda-1)x_0x_1\beta.
\]
Note that $\beta$ is a $(n-2)$-form in $n-1$ variables. Therefore, there exists a polynomial vector field $u = \sum_{i=2}^n g_i \frac{\partial }{\partial x_i}$ such that $g_i = g_1(x_2, \dots,x_n)$ is homogeneous of degree $d-1$ and 
\[
\beta = \iota_u dx_2 \wedge \dots \wedge dx_n.
\]
It follows that $\iota_u \omega = 0$ and $[u,v] =0$. Thus $\tF \cong \opn \oplus \opn(2-d)$ spanned by $v$ and $u$, which is a contradiction. 

{
\medskip\noindent
\textbf{Case 2:} If  $v = (x_0+ x_1)\del{ x_0} + x_1\del{ x_1}$, we have that $ x_0\del{ x_0} + x_1\del{ x_1} ,x_1\del{ x_0} \in \afix(\sF)$. Then $\omega$ is homogeneous with respect to $\{x_0,x_1\}$  and $\lieder_{x_1\del{ x_0}}\omega = 0$, which, combined with $\iota_{\rad}\omega=0$, imply that
\[
\omega = (x_1dx_0-x_0dx_1)\wedge \eta + dx_0\wedge dx_1 \wedge \iota_{\rad}\eta + \gamma
\]
where $\eta$ and $\gamma$ do not depend on $x_0, dx_0,dx_1$. By homogeneity, $\eta$ does not depend on $x_1$ and $\gamma = x_1^2\gamma'$. Then
\[
0 = \iota_v\omega = x_1^2\eta + ((x_0+x_1)dx_1 -x_1dx_0)\wedge \iota_v\iota_{\rad}\eta, 
\]
and we conclude that $\eta = 0$. Therefore, $\omega$ does not depend on $x_0,x_1,dx_0,dx_1$. This implies $h^0(\tF(-1)) >0$ which is a contradiction.


\medskip\noindent
\textbf{Case 3:} If $v = x_1\del{ x_0} + x_2\del{ x_2}$, we have that $ x_2\del{ x_2} ,x_1\del{ x_0} \in \afix(\sF)$. Then, again, 
\[
\omega = (x_1dx_0-x_0dx_1)\wedge \eta + dx_0\wedge dx_1 \wedge \iota_{\rad}\eta +\gamma
\]
where $\eta$ and $\gamma$ do not depend on $x_0, dx_0,dx_1$. And $\omega$ is homogeneous with respect to $x_2$. Then
\begin{align*}
\iota_v\omega & = x_1^2\eta -x_2 (x_1dx_0-x_0dx_1)\wedge \iota_{\del{ x_2}}\eta + x_1dx_1\wedge  \iota_{\rad}\eta + x_2 dx_0\wedge dx_1 \wedge \iota_{\del{ x_2}}\iota_{\rad}\eta + x_2 \iota_{\del{ x_2}}\gamma \\
&= dx_0\wedge ( -x_1 x_2 \iota_{\del{ x_2}}\eta) + 
    dx_1 \wedge( x_0x_2\iota_{\del{ x_2}}\eta + x_1\iota_{\rad}\eta ) +  
        dx_0\wedge dx_1\wedge(x_2 \iota_{\del{ x_2}}\iota_{\rad}\eta) + \\ & \quad
            + x_1^2\eta + x_2 \iota_{\del{ x_2}}\gamma   = 0 
\end{align*}
Therefore, $\iota_{\del{ x_2}}\eta =  \iota_{\rad}\eta =  x_1^2 \eta + x_2\iota_{\del{ x_2}}\gamma = 0$ and
\[
\omega = (x_1x_2dx_0-x_0x_2dx_1-x_1^2 dx_2)\wedge \theta + x_2\rho
\]
where $\eta = x_2\theta$ and $\gamma = -x_1^2dx_2\wedge \theta + x_2\rho$. Moreover, $\theta$ and $\rho$ do not depend on $x_0,x_2,dx_0,dx_1,dx_2$. Note that $\iota_{u}\omega = 0$ for $u = x_0\del{ x_0} + x_1\del{ x_1}$, which commutes with $v$. Therefore, $\tF \cong \opn^{\oplus 2}$, which is a contradiction. 
}
\end{proof}

 

In light of the previous proposition, we need to understand foliations tangent to algebraic actions with codimension-two fixed points. We start with additive $\CC$-actions, which correspond to cases 3) and 4) in Lemma \ref{lem:veccd2} up to a linear change of coordinates. 


\begin{proposition}\label{prop:class-om-nilp}
Let $\sF$ be a foliation of dimension $2$ and degree $2$ on $\pn$ tangent to a global $\CC$-action with codimension-two fixed points. Then $\sF \in \logg{1^{n-2},2}{2}{n}$.
\end{proposition}


\begin{proof}
We fix an action tangent to a vector field in either 3) or 4) in Lemma \ref{lem:veccd2}. Both cases follow the same argument, so we only prove one. In case 3), $v = x_1\del{ x_0} + x_3\del{ x_2}$. Then it has a rational first integral $\rho\colon \pn \dashrightarrow \PP(1^{n-2},2)$ given by 
\[
\rho(x_0: \dots : x_n) = (x_1:x_3:x_4:\dots :x_n: x_0x_3-x_1x_2).
\]
Then $\sF = \rho^*\sH$ for $\sH$ a foliation by curves on $\PP(1^{n-2},2)$. Note that $\rho$ does not ramify in codimension one, hence 
\[
\opn(n+1) = \det(\nF) = \rho^*\det(\nH) = \opn(n+\deg \sH).
\]
Thus $\deg(\sH) = 1$ and $\sH$ is defined by a global vector field. As we observed in Remark \ref{rem:log(1,2)}, $\sF \in \logg{1^{n-2},2}{2}{n}$. 
\end{proof}


Actions of the multiplicative group $\CC^\ast$ with codimension-two fixed locus are,  up to a linear change of coordinates, induced by vector fields as in 1) in Lemma \ref{lem:veccd2}, $v = x_0\del{ x_0} + \mu x_1\del{ x_1}\in H^0(\tpn)$ with $\mu \in \QQ$. The next result classifies the $2$-dimensional foliations invariant under such actions. Note that we don't impose any restriction on $\mu\in \CC^*$ a priori. 

\begin{proposition}\label{prop:class-om-diag}
 Let $\omega\in H^0(\Omega_{\pn}^{n-2}(n+1))$ not vanishing in codimension one and let $v = x_0\del{ x_0} + \mu x_1\del{ x_1}\in H^0(\tpn)$ be such that $\iota_v \omega = 0$ and $\iota_vd\omega = a \omega$ for $a\in\CC$ a constant. Then $\omega$ defines a foliation $\sF$ of one of the following types:
 \begin{enumerate}
 \item $a = \mu + 1$, $\mu \neq -1$ and $\sF \in \logg{1^{n+1}}{2}{n}$;
 \item $a = 0$, $\mu = -1$ and $\sF\in \logg{1^{n-1},2}{2}{n}$;
 \item $a = 0$, $\mu = -2$ and $\sF \in \rat{1^{n-2},3}{n}$;
 \item $\mu \in \{1,2\}$ and $a = \mu + 2$, and $\sF$ is induced by an action of $\mf{aff}(\CC)$.
\end{enumerate}
\end{proposition}


\begin{proof}
 Since $\iota_v\omega = 0$ we may write  
 \[
 \omega = (x_0 dx_1 - \mu x_1 dx_0) \wedge \iota_{\rad}\beta + (\mu-1)x_0x_1\beta \, + \, \gamma,
 \]
 where $\beta$ and $\gamma$ do not depend on $dx_0$ nor $dx_1$ and $\iota_{\rad}\gamma = 0$. Then $\beta$ and $\gamma$ are $(n-2)$-forms depending on $n-1$ basic differentials, which means that there exists a degree-two homogeneous polynomial $q$ and a linear vector field $u = \sum_{j=2}^n b_j \del{ x_j}$ such that $\beta = \iota_u \Theta$ and $\gamma = q \, \iota_{\rad}\Theta$, where $\Theta = dx_2 \wedge \dots \wedge dx_n$. Hence,
 \[
 \omega = (x_0 dx_1 - \mu x_1 dx_0) \wedge \iota_{\rad}\iota_{u} \Theta + (\mu-1)x_0x_1\iota_u \Theta \, + \,q \, \iota_{\rad}\Theta.
 \]
 We might assume that $u$ is not proportional to the radial in the last $n-1$ variables, i.e., $u\wedge \sum_{j=2}^n x_j \del{ x_j} \neq 0$, otherwise $\omega$ would vanish in codimension one. 
 
 By hypothesis, there exists $a\in \CC$ such that $\iota_v d\omega = a \omega$. Computing the LHS, we will see that there are few possibilities for $a$. Start with expanding 
 \begin{align*}
 d\omega & = (\mu+1)dx_0 \wedge dx_1 \wedge \iota_{\rad}\iota_u \Theta - (x_0 dx_1 - \mu x_1 dx_0) \wedge d\iota_{\rad}\iota_u \Theta +\\ & \quad + (\mu-1)(x_0dx_1+x_1dx_0)\wedge \iota_u \Theta + (\mu-1)x_0x_1d\iota_u \Theta + dq\wedge \iota_{\rad}\Theta + \, q \, d\iota_{\rad}\Theta 
 \end{align*}
 then, 
 \begin{align*}
 \iota_vd\omega & = (\mu+1)(x_0 dx_1 - \mu x_1 dx_0)\wedge \iota_{\rad}\iota_u \Theta + (x_0 dx_1 - \mu x_1 dx_0) \wedge \iota_v d\iota_{\rad}\iota_u \Theta + \\ & \quad + (\mu-1)(\mu+1)x_0x_1\iota_u \Theta +(\mu-1)x_0x_1 \iota_v d\iota_u \Theta + v(q)\iota_{\rad}\Theta \\
 & = (x_0 dx_1 - \mu x_1 dx_0)\wedge ((\mu+1)\iota_{\rad}\iota_u \Theta + \iota_v d\iota_{\rad}\iota_u \Theta)+ \\ & \quad + (\mu-1)x_0x_1( (\mu+1)\iota_u \Theta + \iota_v d\iota_u \Theta) + v(q)\iota_{\rad}\Theta 
 \end{align*}
 and by $\iota_v d\omega = a \omega$ imposing we get that 
 \begin{align*}
 \iota_v d\iota_{\rad} \iota_u \Theta &= (a-\mu-1)\iota_{\rad} \iota_u \Theta, \\
 \iota_v d\iota_u \Theta &= (a-\mu-1) \iota_u \Theta + c\, \iota_{\rad} \Theta \\
 v(q) &= a \, q -(\mu-1)c \, x_0x_1
 \end{align*}
 where $c\in \CC$ and the second equation follows from the first. 
 
 Notice that $\iota_{\rad}\iota_u \Theta = \sum_{i,j} p_{ij} dx_2\wedge \dots \wedge \widehat{dx_i} \wedge \widehat{dx_j} \wedge \dots \wedge dx_n $, where each $p_{ij}$ is a quadratic polynomial not having monomials $x_0^2$, $x_0x_1$ nor $x_1^2$. Moreover,
 \[
 \iota_v d\iota_{\rad} \iota_u \Theta = \sum_{i,j} v( p_{ij} ) dx_2\wedge \dots \wedge \widehat{dx_i} \wedge \widehat{dx_j} \wedge \dots \wedge dx_n
 \]
 and we only need to compute the action of $v$ on degree two polynomials. From $v(x_0) = x_0$, $v(x_1) = \mu x_1$ and $v(x_j) = 0$ for $j\geq 2$ we get that 
 \begin{equation} \label{eq:v-action-quad}
 \frac{v(x_ix_j)}{x_ix_j} = \begin{cases} 2, & i=j = 0\\ 2\mu , & i=j=1\\
 (\mu+1) , &i = 0, \, j = 1 \\ 
 1, & i = 0, \, j\geq 2 \\ \mu ,& i = 1,\, j \geq 2 \\ 0 , & i,j\geq 2 
 \end{cases} 
 \end{equation}
 Then, to have $\iota_v d\iota_{\rad}\iota_u \Theta= (a - \mu - 1)\iota_{\rad}\iota_u \Theta$, we need $a-\mu-1\in \{ 0, 1, \mu \}$. 
\\
 \paragraph{\textbf{Case 1.0:}} Suppose that $a-\mu-1 = 0$ and that $\mu \not\in \{ -1,1\}$. Then $\iota_v d \iota_u \Theta = c \, \iota_{\rad} \Theta$ implies $c =0$. Thus $v(q) = (\mu+1)q$. From \eqref{eq:v-action-quad} and $\mu \not\in \{ -1,1\}$, $q = c' x_0x_1$ for $c'\in \CC$ a constant.
 Then 
 \[
 \omega = (x_0 dx_1 - \mu x_1 dx_0) \wedge \iota_{\rad}\iota_u \Theta + x_0x_1((\mu-1)\iota_u \Theta \, + c'\, \iota_{\rad}\Theta). 
 \]
 Consider $w = c'\, x_1 \del{ x_1} + u$ and notice that 
 \[
 \iota_w\iota_v\iota_{\rad} (dx_0\wedge \dots \wedge dx_n) = \iota_w\iota_v\iota_{\rad} (dx_0\wedge dx_1 \wedge \Theta) = \omega. 
 \]
 Also, since $u$ does not depend on $x_0$ nor $x_1$, $[v,w] =0$. Thus $\omega$ defines a foliation $\sF$ on $\pn$ with trivial tangent sheaf $\tF = \opn^{\oplus 2}$. 

 \paragraph{\textbf{Case 1.1:}} Suppose that $a-\mu-1 = 0$ and that $\mu =1$. As in the previous case, $c=0$ and $u$ does not depend on $x_0$ nor $x_1$. But now $v(q) = 2 q$ implies that $q = q_{0}x_0^2 + 2q_{1}x_0x_1+ q_{2}x_1^2$. Then
 \[
 \omega = (x_0 dx_1 - x_1 dx_0) \wedge \iota_{\rad}\iota_u \Theta + (q_{0}x_0^2 + 2q_{1}x_0x_1+ q_{2}x_1^2) \iota_{\rad}\Theta.
 \]
 Consider the vector field
 \[
 w = \left( -q_1x_0 -q_2 x_1 \right)\del{ x_0} + \left(q_0x_0 +q_1x_1 \right)\del{ x_1} + u.
 \]
 As in the previous case, we have that $ \omega = \iota_w \iota_v\iota_{\rad} (dx_0\wedge dx_1 \wedge \Theta) $, and it defines a foliation whose tangent sheaf is trivial.
\\
 \paragraph{\textbf{Case 1.2:}} Suppose that $a-\mu-1 = 0$ and that $\mu =-1$. Again, we have $c=0$ but now $v(q) = 0$, which implies that $q = q_0 x_0x_1 + q'(x_2, \dots, x_n)$. Then
 \[
 \omega = d(x_0x_1) \wedge \iota_{\rad}\iota_u \Theta -2x_0x_1\iota_u \Theta \, + \,q \, \iota_{\rad}\Theta.
 \]
 Then $\omega$ defines a foliation on $\pn$ which is the pullback of the foliation on $\PP(2,1^{(n-1)})$ given by the vector field
 \[
 m = (q_0 y_1 + q'(y_2, \dots, y_n))\del{ y_1} + \sum_{j=2}^n b_j(y)\del{ y_j}
 \]
 via the map defined by $y_1 = x_0x_1$ and $y_j = x_j$ for $j\geq 2$; recall that $u = \sum_{j=2}^n b_j(x)\del{ x_j}$.
 \\
 \paragraph{\textbf{Case 2.0:}} Suppose that $a-\mu-1 = 1$ and that $\mu \not\in \{ -2,-1,2\}$. Then $\iota_v\iota_{\rad}\iota_u \Theta = \iota_{\rad}\iota_u \Theta$ and $x_0$ must divide $\iota_{\rad}\iota_u \Theta$. Hence, there exist $\xi$ and $c'$ such that $\iota_u \Theta = x_0 \xi + c'\iota_{\rad}\Theta$. Note that $c' = -c$. Then
 \[
 \omega = x_0(x_0 dx_1 - \mu x_1 dx_0) \wedge \iota_{\rad}\xi + (\mu-1)x_0^2x_1 \xi + (-c(\mu-1)x_0x_1 + q ) \iota_{\rad}\Theta.
 \]
 On the other hand, $v(q) = (\mu+2)q -(\mu-1) c\, x_0x_1$ implies that $q = c(\mu-1)\, x_0x_1$. Hence, $\omega$ vanishes along $V(x_0)$ which contradicts our hypothesis. 
\\
 \paragraph{\textbf{Case 2.1:}} Suppose that $a-\mu-1 = 1$ and that $\mu = -2$. Then, $v(q) = 3 c x_0x_1$ which implies that $q = -3cx_0x_1 + q'(x_2, \dots, x_n)$. Hence,
 \begin{align*}
 \omega &= x_0(x_0 dx_1 +2 x_1 dx_0) \wedge \iota_{\rad}\xi -3x_0^2x_1 \xi + (3c\,x_0x_1 + q ) \iota_{\rad}\Theta = \\
 & = x_0 [(x_0 dx_1 +2 x_1 dx_0) \wedge \iota_{\rad}\xi -3x_0x_1\xi] + q' \iota_{\rad} \Theta. 
 \end{align*}
 Note that $\xi$ is a constant $(n-2)$-from on $\{dx_2, \dots, dx_n\}$. Then, up to a linear change on the last $n-1$ coordinates, we may assume that $\xi = dx_3\wedge \dots \wedge dx_n$. Then
 \[
 \omega = \iota_{\rad}( (q'dx_2- d(x_0^2x_1))\wedge dx_3 \wedge \dots \wedge dx_n).
 \]
 Now consider $f(x_2, \dots, x_n)$ a degree $3$ polynomial such that $\frac{\partial f}{\partial x_2} = q'$. It follows that 
 \[
 \omega = \iota_{\rad}( d(f-x_0^2x_1)\wedge dx_3 \wedge \dots \wedge dx_n),
 \]
 and $\omega$ defines a rational foliation of type $(3,1^{(n-2)})$. 
 \\
 \paragraph{\textbf{Case 2.2:}} Suppose that $a-\mu-1 = 1$ and that $\mu = 2$. Then, $v(q) = 4q -cx_0x_1$ which implies $q = q_0x_1^2 +c\,x_0x_1$. Then
 \begin{align*}
 \omega &= x_0(x_0 dx_1 - 2x_1 dx_0) \wedge \iota_{\rad}\xi + x_0^2x_1 \xi + (-cx_0x_1 + q ) \iota_{\rad}\Theta\\
 & = x_0(x_0 dx_1 - 2x_1 dx_0) \wedge \iota_{\rad}\xi + x_0^2x_1 \xi + q_0x_1^2 \iota_{\rad}\Theta
 \end{align*}
 As in the previous case, we may assume $\xi = dx_3\wedge \dots\wedge dx_n$, hence 
 \begin{align*}
 \omega &= x_0(x_0 dx_1 - 2x_1 dx_0) \wedge \iota_{\rad}\xi + x_0^2x_1 \xi + q_0x_1^2(x_2\xi -dx_2\wedge \iota_{\rad}\xi) \\
 & = (x_0(x_0 dx_1 - 2x_1 dx_0) - q_0x_1^2 dx_2)\wedge \iota_{\rad}\xi + (x_0^2x_1 + q_0x_1^2x_2)\xi \\
 & = \iota_{\rad}( ( x_0^2 dx_1 - 2x_1x_0 dx_0 - q_0x_1^2 dx_2) \wedge dx_3 \wedge \dots \wedge dx_n).
 \end{align*}
 Note that $q_0 \neq 0$ unless $\omega$ vanishes along $V(x_0)$. Define the vector field $w = \frac{q_0}{2}x_1 \del{ x_0} + x_0 \del{ x_2}$. It follows that 
 \[
 \omega = \iota_{\rad}\iota_w \iota_v ( ( x_0^2 dx_1 - 2x_1x_0 dx_0 - q_0x_1^2 dx_2) \wedge dx_3 \wedge \dots \wedge dx_n).
 \]
 Also notice that $[v,w] = w$ hence $\omega$ defines a foliation associated with an algebraic action of the affine Lie algebra.
\\
 \paragraph{\textbf{Case 2.3:}} Suppose that $a-\mu-1 = 1$ and that $\mu = -1$. Again $x_0$ divides $\iota_{\rad}\iota_u \Theta$. But now $v(q) = q + 2c\, x_0x_1$, which implies that $x_0$ divides $q$. Then $\omega$ vanishes along $V(x_0)$ which is absurd. 
 \\
 \paragraph{\textbf{Case 2.4:}}
 Suppose that $a-\mu-1 = 1$ and that $\mu = 1$. Then $v(q) = 3q$ implies $q=0$. Then, 
 \[
 \omega = \iota_{\rad}\iota_v\iota_u (dx_0\wedge \dots \wedge dx_n). 
 \]
 On the other hand, $\iota_vd\iota_{\rad}\iota_u\Theta = \iota_{\rad}\iota_u\Theta$ implies that $u = u' + c'\sum_{j= 2}x_j \del{ x_j}$, where the coefficients of $u'$ depend uniquely on $x_0$ and $x_1$. Note that we can replace $u$ with $u'$ in the definition of $\omega$. Moreover, $[v,u']=u'$ and $\omega$ defines a foliation tangent to an action of $\mathfrak{aff}(\CC)$.
 \\
 \paragraph{\textbf{Case 3:}} Suppose that $a-\mu-1 = \mu$. We may exchange $x_0$ and $x_1$ so that we get new parameters $(\mu',a')= (1/\mu, a/\mu)$. For $a = 2\mu + 1$ we get $a' = \mu'+2$. Then we reduce to one of the previous cases. 
\end{proof}







% {\color{orange}
% Note that $\lieder_{\del{ x_0}} fdx_J = \frac{\partial f }{\partial x_0}dx_J$. Then $\lieder_{x_1\del{ x_0}} fdx_J = x_1\frac{\partial f }{\partial x_0} + fdx_1\wedge \iota_{\frac{\partial  }{\partial x_0}}dx_J$. Fix
% \[
% \omega = \sum_{|J|  = k} a_J(x) dx_J.
% \]
% Then 
% \[
% \begin{split}
%      0 = \lieder_{x_1\del{ x_0}} \omega = \sum_{\substack{|I|=k-1 \\  0,1\not\in I}} \left( x_1\frac{\partial a_{\{0\}\cup I}}{\partial x_0}dx_0+ \left(a_{\{0\}\cup I} + x_1\frac{\partial a_{\{1\}\cup I}}{\partial x_0}\right)dx_1 \right)\wedge dx_I + \\
%     \sum_{\substack{|I|=k-2 \\ 0,1\not\in I}}x_1\frac{\partial a_{\{0,1\}\cup I}}{\partial x_0}dx_0\wedge dx_1 \wedge dx_I + 
%     \sum_{\substack{|J|=k\\  0,1\not\in J}}x_1\frac{\partial a_J}{\partial x_0}\wedge dx_J 
% \end{split}
% \]
% Therefore, $\frac{\partial a_J}{\partial x_0} = 0$ if $1 \not\in J$ or $0,1\in J$, and $a_{\{0\}\cup I} = -x_1\frac{\partial a_{\{1\}\cup I}}{\partial x_0}$ if $1,0\not\in I$, and $\frac{\partial^2 a_{\{1\}\cup I}}{\partial x_0^2}=0$. We can write $a_{\{1\}\cup I} = c_{\{1\}\cup I} +x_0b_{\{1\}\cup I}$, where $\frac{\partial b_{\{1\}\cup I}}{\partial x_0}= \frac{\partial c_{\{1\}\cup I}}{\partial x_0}=0$. We get that 
% \begin{align*}
%     \omega &= \sum_{\substack{|I|=k-1 \\  0,1\not\in I}} \left( a_{\{1\}\cup I} dx_1 - x_1\frac{\partial a_{\{1\}\cup I}}{\partial x_0}dx_0 \right)\wedge dx_I + dx_0\wedge dx_1 \wedge \theta + \gamma \\
%     & = \sum_{\substack{|I|=k-1 \\  0,1\not\in I}} \left(c_{I}dx_1 +b_{I}(x_0dx_1 - x_1dx_0) \right)\wedge dx_I + dx_0\wedge dx_1 \wedge \theta + \gamma
% \end{align*}
% where $\theta$ and $\gamma$ do not depend on $x_0$. Next, we contract with the radial.
% \[
% \begin{split}
%    0 = \iota_{\rad}\omega = \sum_{\substack{|I|=k-1 \\  0,1\not\in I}} c_{I}x_1 dx_I - \sum_{\substack{|I|=k-1 \\  0,1\not\in I}}\left(c_{I}dx_1 +b_{I}(x_0dx_1 - x_1dx_0) \right)\wedge \iota_{\rad} dx_I \\ 
%    + (x_0dx_1-x_1dx_0)\wedge \theta + dx_0\wedge dx_1 \wedge \iota_{\rad}\theta + \iota_{\rad}\gamma
% \end{split}
% \]
% Then 
% \begin{align*}
%     & \iota_{\rad} \theta =0 \\
%     & \iota_{\rad}\gamma = - \sum_{\substack{|I|=k-1 \\  0,1\not\in I}} c_{I}x_1 dx_I \\
%     &  \sum_{\substack{|I|=k-1 \\  0,1\not\in I}}\left(c_{I}dx_1 +b_{I}(x_0dx_1 - x_1dx_0) \right)\wedge \iota_{\rad} dx_I =(x_0dx_1-x_1dx_0)\wedge \theta
% \end{align*}
% Contracting again the last equation, we get $c_{I}=0$. Hence
% \begin{align*}
%     & \iota_{\rad}\gamma = 0 \\
%     & \theta = \iota_{\rad} \sum_{\substack{|I|=k-1 \\  0,1\not\in I}}b_I dx_I  \\
%     & \eta := \sum_{\substack{|I|=k-1 \\  0,1\not\in I}}b_I dx_I      
% \end{align*}
% Then 
% \[
% \omega = (x_1dx_0-x_0dx_1)\wedge \eta + dx_0\wedge dx_1 \wedge \iota_{\rad}\eta + \gamma
% \]
% where $\eta$ and $\gamma$ do not depend on $x_0$. \textbf{This computation can be shortened. }
% }


% {\color{gray} TO BE OMMITED!
% \begin{lemma}\label{lem:expr-om-ivar}
%  Let $v= p(x)\del{ x_i} + q(x)\del{ x_j}$, with $\gcd(p,q) = 1$, be a homogeneous vector field and let $\omega$ be a homogeneous $p$-form such that $\iota_v \omega =0$ and $\iota_{\rad} \omega = 0$. Then we may write
%  \[
% \omega = (p(x)dx_j - q(x)dx_i)\wedge \iota_{\rad} \beta + (qx_i - p x_j )\beta \, + \, \gamma, 
% \]
% where $\beta$ and $\gamma$ are forms that do not depend on $dx_i$ nor $dx_j$.
% \end{lemma}

% \begin{proof}
% Let us start by writing
% \[
% \theta = dx_i\wedge dx_j \wedge \omega_0 + dx_i\wedge \omega_1 + dx_j\wedge \omega_2 + \omega_3
% \]
% where $\omega_j$ do not depend on $dx_i$ nor $dx_j$. Then, 
% \[
% 0 = \iota_v \omega = (p(x)dx_j - q(x)dx_i)\wedge \omega_0 + p(x)\omega_1 + q(x) \omega_2 
% \]
% which implies $\omega_0 = 0$ and, since $\gcd(p,q) = 1$, there exists $\theta$ such that $\omega_1 = -q(x)\theta$ and $\omega_2 = p(x)\theta$. Hence,
% \[
% \omega = (p(x)dx_j - q(x)dx_i)\wedge \theta + \omega_3.
% \]
% If moreover, we assume that 
% \[
% 0 = \iota_{\rad} \omega = (p(x)x_j-q(x)x_i)\theta - (p(x)dx_j - q(x)dx_i)\wedge \iota_{\rad} \theta + \iota_{\rad} \omega_3
% \]
% we get $\iota_{\rad} \theta = 0$ and $\iota_{\rad} \omega_3 = - (p(x)x_j-q(x)x_i)\theta$. Thus, there exist $\beta$ and $\gamma$ such that $\omega_3 = (q(x)x_i-p(x)x_j)\beta + \gamma$ and $\iota_{\rad}\gamma = 0$. Hence,
% \[
% \omega = (p(x)dx_j - q(x)dx_i)\wedge \iota_{\rad} \iota_u \Theta + (qx_i - p x_j )\beta \, + \, \gamma .
% \]
% \end{proof}
% }



\subsection{3-dimensional foliations with nonstable tangent sheaf.}

 
We focus on the case $c_1(\tF) = 0$, i.e., degree three. Note that, in this case, $\tF$ is stable if and only if $h^0(\tF)= h^0(\tF^\vee) = 0$.  

\begin{proposition}
Let $\sF$ be a 3-dimensional foliation on $\pn$ of degree $3$. Suppose that $\tF$ is not semistable, then either
\begin{itemize}
    \item $\sF$ is a linear pullback of a degree-3 foliation on $\p{n-1}$ or $\p{n-2}$, or
    \item $\sF$ is a pullback of a degree-3 foliation on $\PP(1^{n-2},2)$.
\end{itemize}
\end{proposition}

\begin{proof}
Since $\tF$ is not semistable, either $h^0(\tF(-1)) \neq 0$ or there exists $G\subset \tF$ such that $c_1(G) >0$. The latter is equivalent to $h^0(\tF^\vee(-1)) \neq 0$. In the first case, $\sF$ is a linear pullback of a degree-3 foliation on $\p{n-h}$ for $h = h^0(\tF(-1)) \leq 2$. This gives our first alternative.

Assume $h^0(\tF(-1)) = 0$ and $h^0(\tF^\vee(-1)) \neq 0$. This corresponds to an exact sequence
\[
0 \lra \tG \lra \tF \lra \sI_W(-1) \lra 0 
\]
for some subdistribution $\sG$ and some subscheme $W$. Note that $c_1(\tG) = 1$ and $h^0(\tG(-1)) = 0$. Then, $\tG$ is stable and a maximal destabilizing subsheaf for $\tF$. Hence, $\sG \subset \sF$ is a subfoliation. By the classification of degree-one foliations \cite[Theorem 6.2]{LPT2}, $\sG$ is defined by a rational map $\pn \dashrightarrow \PP(1^{n-2},2)$ given by $n-2$ linear polynomials and one quadratic. This proves the second alternative.
\end{proof}

\begin{proposition}
Let $\sF$ be a 3-dimensional foliation on $\pn$ of degree $3$. Suppose $\tF$ is strictly semistable and indecomposable, i.e., not a direct sum of two other sheaves. Then either
\begin{enumerate}
    \item $\sF$ is tangent to a global vector field singular in codimension at most 3;
    \item $\sF$ is tangent to a 2-dimensional Lie algebra action singular in codimension 2;
    \item $\sF$ admits a 2-dimensional subfoliation $\sG$ of degree 2  such that $\tG$ is stable;
    \item $\tF$ is defined by a nontrivial extension $0 \to \opn \to \tF^\vee \to Q$, where $Q$ is a stable reflexive sheaf. 
\end{enumerate}

\end{proposition}

\begin{proof}
First, $\tF$ being strictly semistable means that $h^0(\tF(-1)) = h^0(\tF^\vee(-1)) =0 $ and either $h^0(\tF) \neq 0$ or $h^0(\tF^\vee)\neq 0$.  By Lemma \ref{lem:sstab-h0}, $h^0(\tF) \leq 2$; otherwise $\tF \cong \opn^{\oplus 3}$. If $h^0(\tF) = 1$, then we have a global section $\sigma$ and the corresponding vector field $v_\sigma$. If $v_\sigma$ vanishes in codimension at least $4$, then $\tF = \opn \oplus \tG$, by de Rham-Saito division, see \cite[\S1]{livro-alcides}. Hence, $v_\sigma$ must vanish in codimension at most $3$. This gives the first alternative.

Suppose $h^0(\tF) = 2$. By Proposition \ref{prop:sstable-subfol}, $\sF$ admits a subfoliation $\sG$ such that $\tG = \opn^{\oplus 2}$. By Proposition \ref{prop:sing2xsplit}, $\tF/\tF = \sI_W$ for $\emptyset \neq W \subset \sing(\sG)$ of codimension two. This gives the second alternative.

Next, suppose $h^0(\tF) = 0$. Then $h^0(\tF^\vee)= 1$. Indeed, if $h^0(\tF^\vee) \geq 3$ then $\tF = \opn ^{\oplus 3}$, and if $h^0(\tF^\vee) =2$ then $h^0(\tF) \geq 1$. Let $\sigma \in H^0(\tF^\vee)$. We get the following exact sequence. 
\begin{equation}\label{eq:ext00}
    0 \lra \tG \lra \tF \stackrel{\sigma^\vee}{\lra} \sI_W \lra 0
\end{equation}
where $\tG$ is stable of rank two, and $W$ is the degeneracy locus of $\sigma$. If $\sG$ is integrable, we get the third alternative. Otherwise, the Lie bracket induces a nonzero morphism $\extp^2\tG \to \sI_W$. Note that this morphism kills the torsion part $T$ of $\extp^2\tG$ and $\extp^2\tG/T = \sI_Y$ for $Y$ supported at the singular locus of $\tG$. Taking double duals, one sees that any morphism $\sI_Y \to \sI_W$ is given by multiplication with a constant. We conclude that $W\subset Y$. In particular, $\codim W \geq 3$. Recall that $\inext^1(\sI_W, \opn)$ is supported on the codimension-two part of $W$, which is empty. Therefore, dualizing \eqref{eq:ext00} one gets
\[
0 \lra \opn \lra \tF^\vee \lra \tG^\vee \lra \inext^1(\sI_W, \opn) = 0
\]
proving the last alternative.
\end{proof}

By a result of Druel, a foliation with stable tangent sheaf $c_1=0$ and dimension $\leq 3$ must be algebraically integrable. Thus, in the third alternative, the foliation $\sF$ is a pullback of a foliation by curves. 

An interesting question is whether an analog of Proposition \ref{prop:nonsplit} holds for 3-dimensional foliations such that $h^0(\tF) \geq 2$. Following the same strategy, one is led to compute $\sG$, the subfoliation defined by $\fix(\sF)$. The computations of $\afix(\sF)$ become too cumbersome, and we didn't pursue this investigation.










%\section{Foliations of dimension 2 and degree 2 }\label{degree2-dim2}
\section{Main result}\label{sec:main}
In this section, we prove Theorem \ref{thm:mainA}. We start with the case of two-dimensional foliations. 

\begin{theorem}\label{thm:dim2-structure}
Let $\sF$ be a two-dimensional foliation on $\pn$ of degree two. Then one of the following holds:
\begin{enumerate}
    \item $\tF$ is stable and $\sF$ is algebraically integrable with uniruled leaves;
    \item $\tF$ is strictly semistable and $\sF$ belongs to either 
        \begin{itemize}
            \item an Affine-Lie component $\exc{d_1, \dots, d_r}{2}{n}$;
            \item one of the rational components $\rat{1^{n-2},2,2}{n}$ or $\rat{1^{n-1},3}{n}$; or
            \item one of the logarithmic components $\logg{1^{n-1},2}{2}{n}$ or $\logg{1^{n+1}}{2}{n}$.
        \end{itemize}
    \item $\tF$ is not semistable and $\sF$ is a linear pullback of a degree-$2$ foliation by curves on $\p{n-1}$.
\end{enumerate}
\end{theorem}
 
\begin{proof}
Given a foliation $\sF$ as in the statement of the theorem, its tangent sheaf $\tF$ is reflexive of rank two with zero first Chern class. Then $\tF$ is semistable if and only if $h^0(\tF(-1)) = 0$, and stable if and only if $h^0(\tF) = 0$.  If $\tF$ is not semistable, then $\tF \cong \opn(1)\oplus \opn(-1)$ and $\sF$ is a linear pullback of a degree-$2$ foliation by curves on $\p{n-1}$. This proves item \textit{3)}.


If $\tF$ is stable then $c_2(\tF) > 0$, otherwise it would be trivial, see \cite[Lemma 2.5]{JG}. Since $\pn$ is smooth and simply connected, it follows from Druel's results \cite[Theorem 6.1 and Proposition 8.4]{Druel} that $\sF$ is algebraically integrable. Due to \cite[Theorem 5.13]{LPT}, $\sF$ has uniruled leaves. This proves the first item. 

It remains to understand the case where $\tF$ is strictly semistable. If $\tF = \opn^{\oplus 2}$ then $\sF$ is associated with a Lie algebra action, and thus it falls in either $\logg{1^{n+1}}{2}{n}$ or an Affine-Lie component from Theorem \ref{thm:lie-comps}. If $\tF \not\cong \opn^{\oplus 2}$  then Proposition \ref{prop:sub-deg2-alg} says that $\sF$ is tangent to a one-dimensional algebraic action with codimension-two fixed locus. 

If $\sF$ is tangent to an additive action, due to Proposition \ref{prop:class-om-nilp}, then $\sF \in \logg{1^{n-2},2}{2}{n}$. If $\sF$ is tangent to a multiplicative action, Proposition \ref{prop:class-om-diag} implies that $\sF$ must fit in one of the cases already listed.
\end{proof}







%\section{Foliations of degree two and higher dimension}\label{ProofThm1}
Next, we treat the case of dimension $k>2$. Then, the positivity of $\tF$ simplifies the classification. Indeed,  $c_1(\tF) = k-2$, and $\sF$ is a Mukai foliation, see \cite{AD-mukai}. We have the following theorem (cf. Theorem \ref{thm:mainA}).

\begin{theorem}\label{thm:dim>2}
Let $\sF$ be a degree-$2$ foliation on $\pn$ of dimension $k>2$. Then one of the following holds:
\begin{enumerate}
    \item $\sF$ is algebraically integrable;
    \item $\sF \in \logg{1^{n-k+1},2}{k}{n}$ is a logarithmic foliation of type $(1^{n-r +1}, 2)$;
    \item $\sF \in \logg{1^{n+1}}{k}{n}$ is a logarithmic foliation of type $(1^{n+1})$;
    \item  $\sF \in \exc{r_1, \dots,r_m}{k}{n}$ is an Affine-Lie foliation; or
    \item $\sF \in \lpb{2}{k}{n}$ is a linear pullback of a foliation by curves in $\p{n-k+1}$.
\end{enumerate}
\end{theorem}
The proof of this result will follow from Proposition \ref{prop:ra=r-1=log} below. We begin assuming that $\sF$ is not a linear pullback of a foliation on $\p{r}$ with $r<n$. This is equivalent to assuming that $\sF$ has no subfoliation of degree $0$ or that $H^0(\tF(-1)) = 0$. Indeed, for any foliation $\sG$ of degree $0$, we have $\tG \cong \opn(1)^{\oplus \dim \sG}$.

\begin{proposition}\label{prop:ra=r-1=log}
Let $\sF$ a degree-$2$ foliation on $\pn$ of dimension $k >2$. If $H^0(\tF(-1)) = 0$ then either 
\begin{itemize}
    \item $\sF$ is algebraically integrable, or
    \item $\sF \in \logg{1^{n-k+1,2}}{k}{n}$ is a logarithmic foliation of type $(1^{n-r +1}, 2)$. 
\end{itemize}
\end{proposition}

\begin{proof} 
Suppose that $\sF$ is not algebraically integrable. By \cite[Proposition 7.5]{AD} there exists a subfoliation $\sG \subset \sF$ such that $c_1(\tG) \geq c_1(\tF)$, hence, 
\[
\dim \sG - \deg \sG \geq k - 2 > \dim \sG -2
\]
which implies that $\deg \sG \leq 1$. Since $H^0(\tF(-1)) = 0$, we have $\deg \sG = 1$ and $H^0(\tG(-1)) = 0$; hence, $\dim \sG = k-1$. By Theorem \ref{thm:fols-deg1}, $\sG$ is defined by a rational map $\rho \colon \pn \dashrightarrow \PP(1^{n-k+1}, 2)$ given by $n-k+1$ linear forms and one quadratic form. It follows that $\sF = \rho^*\sH$ for some (purely transcendental) foliation by curves $\sH$ on $\PP(1^{n-k+1}, 2)$. Since $\rho$
does not ramify in codimension one, $\det(\nF) = \rho^*\det(\nH)$. Hence, $\sH$ is given by a global vector field. This means that $\sF \in \logg{1^{n-k+1},2}{k}{n}$, see Remark \ref{rem:log(1,2)}. 
\end{proof}

Now, we apply this proposition to prove the general case. 

%%%%---------------------
\begin{proof}[Proof of Theorem \ref{thm:dim>2}]
Suppose that $\sF$ is neither algebraically integrable nor logarithmic of type $(1^{n-k+1},2)$. By Proposition \ref{prop:ra=r-1=log}, we must have $h = h^0(\tF(-1)) >0$. Then $\sF$ is a linear pullback of a foliation $\sG$ on $\p{n-h}$ such that $H^0(\tG(-1)) = 0$. It follows that $\dim \sG \leq 2$. Indeed, if we had $\dim \sG \geq 3$ we would get a contradiction applying Proposition \ref{prop:ra=r-1=log} again. 

If $\dim \sG = 1$ then $\sF \in \lpb{2}{k}{n}$. If $\dim \sG = 2$ then, by Theorem \ref{thm:dim2-structure}, $\sG$ belongs to either $\logg{1^{n-h+1}}{2}{n-h}$ or a Affine-Lie component $\exc{r_1, \dots,r_m}{2}{n-h}$. Thus $\sF$ belongs to either  $\logg{1^{n+1}}{k}{n}$ or $\exc{r_1, \dots,r_m}{k}{n}$, see Remark \ref{rem:exceptional}.
\end{proof}


\section{Holomorphic Poisson structures}\label{sectionPoisson}
In this section, we will prove Corollary \ref{Cor:Poisson}. A holomorphic Poisson structure on a complex manifold $X$ is a $\mathbb{C}$-linear Lie bracket 
$$\{ \cdot , \cdot \}\colon \OO_X \times \OO_X \to \OO_X$$ 
which satisfies the Leibniz rule $\{ f , gh \} = h \{f,g \} + f\{ g,h\}$ 
and Jacobi identities
$$\{f, \{g, h\}\} + \{g, \{h, f\}\} + \{h, \{f, g\}\} = 0$$
for all $f,g \in \OO_X$. The bracket corresponds to a holomorphic bivector field $\sigma \in H^0(X, \wedge^2 TX)$ given by
$\sigma(df \wedge dg) = \{f, g\}$, for all $f,g \in \OO_X$. 
We will denote a Poisson structure on $X$ as the pair $(X, \sigma)$, where $\sigma \in H^0(X, \wedge^2 TX)$ is the corresponding Poisson bivector field. We say that the structure $(X, \sigma)$ has generic rank $2k$ if 
\begin{center}
 $\underbrace{\sigma \wedge \cdots \wedge \sigma}_{k\ \mathrm{times}} \neq 0 \quad \text{and} \quad \underbrace{\sigma \wedge \cdots \wedge \sigma}_{k+1\ \mathrm{times}} = 0$. 
\end{center} 
The holomorphic bivector $\sigma$ induces a morphism
$$
\sigma^{\#}\colon \Omega_X^1 \to TX
$$
which is called the {\it anchor map} and is defined by $\sigma^{\#}(\theta) = \sigma(\theta,\cdot)$, where $\theta$ is a germ of a holomorphic $1$-form, and 
$\delta\colon TX \to \wedge^2 TX$ is given by $\delta(v) = L_v(\sigma)$; here $L_v$ denotes the Lie derivative. We refer to \cite{Pym3} for more details.



\begin{definition}
The {\it symplectic foliation} associated with  $\sigma$ is the foliation given by $\sF_{\sigma}:= \mathrm{Ker}(\sigma^{\#})$, whose dimension is the rank of  the anchor map 
$\sigma^{\#}: \Omega_X^1\to TX$. 
A Poisson variety $(X,\sigma)$ is called \emph{generically symplectic} if the anchor map $\sigma^{\#}\colon \Omega_X^1\to TX$ is generically an isomorphism. 
Then, the degeneracy loci of $\sigma^{\#}$ is an effective anti-canonical divisor $D(\sigma)\in |-K_X|$.
\end{definition}



\begin{definition}
Let $(X,\sigma)$ be a Poisson projective variety. A {\it Poisson connection} on a sheaf of $\OO_X$-modules $E$ is a 
$\mathbb{C}$-linear morphism of sheaves 
$\nabla: E \to TX \otimes E$ satisfying the Leibniz rule
$$
\nabla(fs) = \delta(f) \otimes s + f \nabla(s),
$$
where $f$ is a germ of a holomorphic function on $X$ and $s$ is a germ of a holomorphic section of $E$. We say that $E$ is a \emph{Poisson module} if it admits a Poisson flat connection, i.e., if its curvature $\nabla^2\colon E \to \Omega_X^2 \otimes E$ vanishes.
Equivalently, a Poisson connection defines a $\mathbb{C}$-linear bracket $\{ \ , \ \} \colon \OO_X \times E \to E$ by
$$
\{ f, s \} := \nabla(s)(df), 
$$
where $f$ is a germ of a holomorphic function on $X$ and $s$ is a germ of a holomorphic section of $E$. \end{definition}

Let $E\cong \OO_{X}^{\oplus 2}$ be a Poisson module with a trace free Poisson connection $\nabla$, then we can write 
$$
 \nabla=\delta +\begin{pmatrix}v_1 &v_2
 \\v_0&-v_1\end{pmatrix} 
$$
and the flatness condition is equivalent to 
$$ 
 \begin{matrix}
\delta(v_0)=\hfill v_0\wedge v_1\\
\delta(v_1)=2v_0\wedge v_2\\
\delta(v_2)=\hfill v_1\wedge v_2
\end{matrix} .
$$
We refer to \cite{Morita} for more details. 
 \begin{definition}
 Let $D$ be a reduced divisor on $X$. The \emph{logarithmic tangent sheaf} $TX(-\log(D))$ is the subsheaf of $TX$ consisting of germs of vector fields that preserve the ideal sheaf $\OO_X(-D)$. That is, $TX(-\log(D))$ is the sheaf of vector fields tangent to $D$. A reduced divisor $D \subset X$ is \emph{free} if $TX(-\log(D))$ is locally free. If $f : X \to Y$ is a rational map, we denote by $T_{X/Y} \subset TX$ the relative tangent sheaf and by $T_{X/Y}(-\log(D)) := T_{X/Y} \cap TX(-\log(D))$.
\end{definition}


\begin{definition}
An analytic subspace $Z\subset (X,\{ \ , \ \})$ is a \emph{Poisson subspace} 
$Z$ is equipped with a Poisson structure $\{\ , \ \}_{Z}$
such that the inclusion $i\colon Z \to X$ is a morphism of Poisson analytic spaces, i.e., it is compatible with the brackets. 
\end{definition}




\begin{proposition}\cite[Proposition 4.4.1]{Pym3} \label{Poisson-divisor}
Let $D\subset X$ a Poisson free
divisor with respect to a holomorphic bi-vector $\sigma \in H^0(X, \wedge^2TX)$. Then $ \sigma \in H^0(X, \wedge^2TX(-\log(D))$. 
\end{proposition}





 
  

\begin{proof}[Proof of Corollary \ref{Cor:Poisson}]
Since the associated symplectic foliation $\sF_\sigma$ has no divisorial zeros and
$c_1(\tF_\sigma) = 0$, $\sF_\sigma$ is a foliation of dimension $2$ and degree $2$
on $\PP^{4}$. It follows from Theorem \ref{thm:dim2-structure} that one of the
following holds:
\begin{enumerate}
\item[(a)] $\sF_\sigma$ is algebraically integrable and $\sigma$ is not decomposable,
since $\tF_\sigma$ is stable.

\item[(b)] $\sF_\sigma$ is the linear pull-back of a degree-two foliation by curves
on $\PP^{3}$ with tangent sheaf split. So, $\sigma$ is decomposable.

\item[(c)] $\tF_\sigma \cong \mathfrak{g}\otimes \OO_{\PP^{4}}$, where
$\mathfrak{g} \subset \mathfrak{sl}(5,\mathbb{C})$ and either $\mathfrak{g}$ is an
abelian Lie algebra of dimension $2$ or $\mathfrak{g} \cong \mathfrak{aff}(\mathbb{C})$.
So, $\sigma$ is decomposable.

\item[(d)] $\sF_\sigma$ is the pull-back by a rational map
$\rho\colon \PP^{4} \dashrightarrow \mathbb{P}(1^{(3)},2)$ of a non-algebraic foliation
by curves induced by a global vector field on $\mathbb{P}(1^{(3)},2)$. Then $\tF_\sigma$
is not split, which implies that $\sigma$ is not decomposable.
\end{enumerate}

\noindent \textbf{Case (a).}
Since $\sigma$ is of normal crossing type, by Remark \ref{rem:Integraldeg2} the
symplectic foliation $\sF_\sigma$ is tangent to a dominant rational map
$\PP^{4} \dashrightarrow \PP^{2}$ given either by
\[
f_{(1,3)} := (L_1^3, L_2^3, P)
\quad\text{or}\quad
f_{(1,2,2)} := (L_1^2, Q_1, Q_2),
\]
where $L_i$ are linear, $Q_i$ are quadratic, and $P$ is cubic. Then the divisors
\[
D_{(1,3)} := \{L_1 L_2 P = 0\}
\quad\text{and}\quad
D_{(1,2,2)} := \{L_1 Q_1 Q_2 = 0\}
\]
are Poisson divisors of $\sigma$. If we denote by
$T_{\PP^{4}/\PP^{2}}$ the corresponding relative tangent sheaf either of
$f_{(1,3)}\colon \PP^{4} \dashrightarrow \PP^{2}$ or of
$f_{(1,2,2)}\colon \PP^{4} \dashrightarrow \PP^{2}$, then it follows from
Proposition \ref{Poisson-divisor} that
\begin{itemize}
 \item $\sigma\in H^0(\PP^{4}, \wedge^2T_{\PP^{4}/\PP^{2}}(-\log D_{(1,3)}))$, or
 \item $\sigma\in H^0(\PP^{4}, \wedge^2T_{\PP^{4}/\PP^{2}}(-\log D_{(1,2,2)}))$.
\end{itemize}

\medskip
\noindent\textbf{Case (b).}
If $\sF_\sigma$ is the linear pull-back of a degree-two foliation by curves on
$\PP^{3}$, then there are $v_0$ a constant vector field and $v$ a homogeneous
polynomial vector field of degree $2$ on $\PP^{3}$ such that $\sigma = v_0 \wedge v$.
It is clear that $[v_0,v]=0$. Let $H$ be the Poisson hyperplane determined by
$v_0\in \mathbb{P}H^0(\PP^{4}, T \PP^{4}(-1))\cong \check{\PP}^{4}$. That is, $H$ is
the hyperplane parameterizing lines passing through a point which are leaves of the
foliation induced by $v_0$. It is clear that $H$ is a Poisson divisor.

\medskip 
\noindent\textbf{Case (c).}
If $\tF_\sigma \cong \mathfrak{g}\otimes \OO_{\PP^{4}}$, where
$\mathfrak{g}\subset \mathfrak{sl}(5,\mathbb{C})$ is an abelian Lie algebra, then
since $\sF_\sigma$ is general, the Lie algebra $\mathfrak{g}$ is generated by diagonal
linear vector fields which are given, in a suitable system of homogeneous coordinates
$[z_0:\dots:z_4]$ on $\PP^{4}$, by
\[
v_0=\sum_{i=0}^{4} \lambda_{0i}z_{i}\dfrac{\partial}{\partial z_i}
\quad\text{and}\quad
v_1=\sum_{j=0}^{4} \lambda_{1j}z_{j}\dfrac{\partial}{\partial z_j}.
\]
Then
\[
\sigma=v_0\wedge v_1=\sum_{i<j}(\lambda_{0i}-\lambda_{1j})\,z_i z_j\,
\dfrac{\partial}{\partial z_i}\wedge \dfrac{\partial}{\partial z_j}.
\]
Since $v_r(z_i)= \lambda_{ri}z_{i}$ for all $r=0,1$ and $i=0,\dots,4$, we conclude that
$v_0,v_1\in H^0(\PP^{4}, T\PP^{4}(-\log D))$ and
$\sigma\in \wedge^2 \mathfrak{h}$, where
\[
\mathfrak{h}=H^0(\PP^{4}, T\PP^{4}(-\log D))
\]
is the abelian Lie algebra of vector fields tangent to the arrangement of coordinate
hyperplanes
\[
D=\{z_0 z_1 z_2 z_3 z_4=0\}.
\]

Suppose now that $\tF_\sigma \cong \mathfrak{g}\otimes \OO_{\PP^{4}}$, where
$\mathfrak{g}\cong \mathfrak{aff}(\mathbb{C})$. In this case, we have that the
Poisson structure on $\PP^{4}$ is induced by the so-called infinitesimal action
\[
\sigma \colon \mathfrak{g}\otimes \OO_{\PP^{4}} \longrightarrow T\PP^{4},
\]
which is generated by linear vector fields $v_0$ and $v_1$ such that
$[v_0,v_1]=v_0$. In addition, in both cases where $\sigma=v_0\wedge v_1$, with
$v_0, v_1\in H^0(\PP^{4}, T\PP^{4})$, we have that the tangent sheaf of the
symplectic foliation is $\tF_\sigma \cong \OO_{\PP^{4}}^{\oplus 2}$. We have a
trace-free flat Poisson connection 
\[
\nabla\colon \tF_\sigma \longrightarrow T\PP^{4} \otimes \tF_\sigma
\]
given by 
\[
 \nabla =\delta +
 \begin{pmatrix}
   v_1 & 0 \\
   v_0 & -v_1
 \end{pmatrix}.
\]
If $\tF_\sigma \cong \mathfrak{g}\otimes \OO_{\PP^{4}}$ with $\mathfrak{g}$ abelian,
then for $i=0,1$ we have 
\[
 \delta(v_i)= L_{v_i} (v_0\wedge v_1 )= 0,
\]
since $[v_1,v_0]=0$. If $\tF_\sigma \cong \mathfrak{g}\otimes \OO_{\PP^{4}}$ with
$\mathfrak{g}\cong \mathfrak{aff}(\mathbb{C})$, then 
\begin{align*}
    \delta(v_0) &= L_{v_0} (v_0\wedge v_1 )= v_0\wedge v_1,\\
    \delta(v_1) &=L_{v_1} (v_0\wedge v_1 )=0,
\end{align*}
since $[v_1,v_0]=-v_1$. Now, let
$p\in \PP^{4}\setminus Sing(\sF_\sigma)$ and let $L_p$ be the symplectic leaf of
$\sF_\sigma$ passing through $p$. Since
\[
 (\sigma^{\#} )^{-1}\circ \delta |_{L_p}
 \;=\;
 d \circ (\sigma^{\#} )^{-1} |_{L_p},
\]
where $d$ denotes the de~Rham differential, we have that
\[
 \begin{matrix}
 d \circ (\sigma^{\#} )^{-1}(v_0)
 =
 (\sigma^{\#} )^{-1}\circ \delta |_{L_p} (v_0)
 =
 (\sigma^{\#} )^{-1}(v_0)\wedge (\sigma^{\#} )^{-1}( v_1)\\[4pt]
 d \circ (\sigma^{\#} )^{-1}(v_1)
 =
 (\sigma^{\#} )^{-1}\circ \delta |_{L_p} (v_1)
 =
 2 (\sigma^{\#} )^{-1}( v_0)\wedge (\sigma^{\#} )^{-1}(v_2).
\end{matrix} 
\]
Now, defining $(\sigma^{\#} )^{-1} |_{L_p} (v_i)=\omega_i$, we obtain the transversely
affine structure on $L_p$ given by the pair $(\omega_0,\omega_1)$, where $\omega_0$ is
the $1$-form inducing the transversely affine holomorphic foliation on $L_p$. We have:
\begin{itemize}
 \item if $\mathfrak{g}$ is abelian, then $(\omega_0,\omega_1)$ is Euclidean; 
 \item if $\mathfrak{g}\cong \mathfrak{aff}(\mathbb{C})$, then $(\omega_0,\omega_1)$
 is affine. 
\end{itemize}
If $\tF_\sigma \cong \mathfrak{g}\otimes \OO_{\PP^{4}}$, where
$\mathfrak{g}\cong \mathfrak{aff}(\mathbb{C})$, then the normal form of $\sigma$ is
given by the vector fields in \eqref{eq:basic-nilp} specialized to $n=4$.

\medskip 
\noindent\textbf{Case (d).}
The symplectic foliation $\sF_\sigma$ is the pull-back by a rational map
\[
\rho\colon \PP^{4} \dashrightarrow \mathbb{P}(1^{(3)}, 2)
\]
of a foliation by curves on $\mathbb{P}(1^{(3)}, 2)$, and
\[
\rho := (L_1, L_2, L_3, Q),
\]
where $L_i$ are linear and $Q$ is quadratic. Then
\[
\sigma \in H^0(\PP^{4}, \wedge^2T\PP^{4}(-\log D)),
\quad
D := \{L_1 L_2 L_3 Q = 0\}.
\]
Consider
\[
v \in H^0(\PP^{4}, T\PP^{4}(-\log D)) \subset H^0(\PP^{4},T\PP^{4}),
\]
the global vector field whose leaves are fibers of
$\rho\colon \PP^{4} \dashrightarrow \mathbb{P}(1^{(3)}, 2)$. It is clear that
$\sigma \wedge v = 0$ and $dF(v) = 0$, where $F := L_1 L_2 L_3 \cdot Q$. That is,
$v$ is Hamiltonian with respect to $\sigma$.


Finally, in the case $\sigma = v_0 \wedge v_1 \in \wedge^2 \mathfrak{h}$, where
$\mathfrak{h} = H^0(\PP^{4}, T\PP^{4})$ and $v_0$ and $v_1$ are linear
vector fields such that $[v_0,v_1]=v_0$, the normal form follows from subsection 
 \ref{normal-form-poisson}. 
\end{proof}

\begin{remark}
 Let $(\PP^{4}, \sigma)$ be a holomorphic Poisson structure such that
 $\sigma \wedge \sigma = 0$. Then the associated symplectic foliation $\sF_\sigma$
 has degree at most $2$. If $\sigma$ has zeros of codimension one, then by the
 classification of foliations of degree $0$ and $1$ (see Section \ref{sect:folsonPn}),
 we have that either
 \begin{itemize}
   \item $\sigma = Q \cdot (v_0 \wedge v_1)$, where $v_i$ are constant vector fields
   and $Q$ is of degree $2$, or
   \item $\sigma = H_0 \cdot \Tilde{\sigma}$, with $\{H_0=0\}$ being a Poisson
   hyperplane and
   \begin{itemize}
     \item $\Tilde{\sigma} = v_0 \wedge v$, where $v_0$ is a constant vector field,
     $v$ is a linear vector field such that $[v_0, v] = 0$; 
     \item $\Tilde{\sigma} \in H^0(\PP^{4}, \wedge^2 T_{\PP^{4}/\PP^{2}}(-\log D))$,
     where $D = H_1 \cup H_2 \cup H_3 \cup \mathcal{Q}$ is a Poisson divisor such
     that $H_i$ are linear and $\mathcal{Q}$ is of degree $2$. 
   \end{itemize}
 \end{itemize}
\end{remark}
 




 
 

\bibliographystyle{abbrvurl}
\bibliography{bibliography}




\end{document}
